\documentclass{article}
\usepackage[margin=0.8in]{geometry}
\usepackage{amsmath,amssymb,amsthm,mathtools}
\usepackage{mathrsfs,bm,bbm}
\usepackage{graphicx,booktabs,array,multirow,tabularx,longtable}
\usepackage{enumitem}
\usepackage{xcolor}
\usepackage{algorithm,algpseudocode}
\usepackage{subcaption,tikz}
\usepackage{siunitx,cancel,accents}
\usepackage{microtype}
\usepackage[numbers,sort&compress]{natbib}
\usepackage[colorlinks=true,linkcolor=blue,citecolor=blue,urlcolor=blue]{hyperref}
\usepackage{cleveref}
\setlength{\emergencystretch}{2em}
\newtheorem{assumption}{Assumption}
\newtheorem{definition}{Definition}
\newtheorem{theorem}{Theorem}
\newtheorem{lemma}{Lemma}
\newtheorem{proposition}{Proposition}
\newtheorem{openendedquestion}{Open-ended Question}
\newtheorem{conjecture}{Conjecture}
\newcommand{\coreref}[1]{\ref{#1}}

\begin{document}

\sloppy

\section*{scm\_\allowbreak{}hidden\_\allowbreak{}recant\_\allowbreak{}modulus\_\allowbreak{}cert}

\noindent\textit{Specialization:} v1\par

\paragraph{TL;DR.} The principal result identifies the Wu--Bai--Cui same-unit hidden-recanting mean on an interpretable calibrated primitive subclass: binary measured severity, affine treatment-one mediator risk, and additive treatment-zero outcome means directly generate a latent-valid observed bridge, and the observable loading row-space condition makes its scalar value invariant despite failure of three-state proxy completeness. The broad WBC result is retained only as a conditional envelope; its latent-valid bridge intersection automatically supplies observed solvability, so its only additional observable criterion is loading membership in the row space of the six-equation operator. Both results condition on two proposal-specific, non-testable target-relevant latent transport-overlap restrictions, use only the six displayed conditional-moment denominators, and allow zero probabilities for values of \(W\) and \(Y\). An exact finite canonical response-type characterization gives the full causal-fiber range \(\Gamma(P)\), its infimum and supremum, the maximal point-identification subclass by endpoint equality, and checkable two-feasible-point nonidentification certificates. Duarte et al. are credited for the generic discrete response-program and optimization machinery; the characterization here is the WBC bridge-supported specialization and is not a new terminating or certified endpoint solver. Exact rational artifacts give a calibrated positive success SCM with target \(1/2\) and an observationally equivalent failure pair with targets \(19/48\) and \(7/16\), now with an explicit shared-\(U_H\) joint response map. A fixed-channel theorem sharply restores the two WBC completeness clauses with one added binary proxy; WBC identification still requires all remaining premises and bridge existence. Emptiness of the positive-interior obstruction set and root-\(n\) inference remain open.

\section{Project justification}

\paragraph{Gap.} Wu--Bai--Cui use completeness to lift observed bridges to the shared-witness causal target. Zhang--Li--Miao--Tchetgen Tchetgen give abstract functional invariance, while Proximal ID retains ordinary or sequential injectivity and identifies recursive interventional margins rather than this same-unit nested mean. Bai--Wu--Sun--Cui use an observed treatment-induced recanting witness distinct from a latent baseline confounder, three stage-specific completeness conditions, and three outcome bridges to identify a different path-specific mean. Duarte--Finkelstein--Knox--Mummolo--Shpitser provide the generic response-function polynomial reduction, sharp-bound theorem, and anytime optimization machinery for discrete causal problems, but do not provide the fixed WBC bridge-supported class, calibrated below-completeness causal lift, rescue branches, or exact WBC certificates. None of the within-graph comparators supplies a primitive below-completeness subclass for the fixed hidden-witness WBC graph together with full success and failure SCMs.

\paragraph{Niche.} Binary proxies cannot be complete for a three-state hidden witness, yet calibrated common-shock structure can identify the same-unit nested mean at rank-two proxy laws. Beyond that sufficient branch, the finite response-function representation characterizes the exact full causal-fiber range \(\Gamma(P)\), so point identification is characterized by equality of its endpoints rather than by bridge row span alone.

\paragraph{Fill.} The project derives sequential bridge realizability and the causal equality directly in the calibrated subclass, retains a transparent broad conditional envelope, and proves the exact WBC full-law fiber equality \(\Gamma(P)=\{T(\theta):(\theta,x)\in\mathcal F(P)\}\), yielding the maximal point-identification subclass by equality of its infimum and supremum and finite two-point nonidentification certificates. It verifies the exact positive success SCM, completes the observationally equivalent failure pair with a shared-\(U_H\) response map, and proves sharp one-bit restoration of the two completeness clauses for the fixed rank-two channel. The generic response-program and optimization machinery is credited to Duarte et al.; no new terminating solver, endpoint-attainment result, certified approximation, or inference theorem is claimed. Identification after augmentation remains conditional on every other Wu--Bai--Cui premise and observed bridge existence.

\section{Related work}

Wu--Bai--Cui (2026) Assumptions 3.1--3.3, Proposition 3.1, and Theorem 3.1 supply the exact hidden-recanting graph, sequential bridges, completeness-based causal lift, and proximal outcome-regression formula used as the primary within-graph benchmark. Their efficient-influence-function result additionally uses bridge uniqueness and operator surjectivity; neither root-\(n\) inference nor efficiency is claimed here. The calibrated subclass retains the WBC consistency, sequential, cross-world, and proxy-exclusion premises but replaces the use of completeness for bridge validity with direct common-shock algebra. Its two target-relevant latent transport-overlap restrictions are additional proposal-specific cross-world restrictions, not standard WBC positivity: they are not testable from \(O\), depend on design or subject-matter knowledge, and are conditioned on rather than derived by the calibrated construction. Bai--Wu--Sun--Cui (2026), Section 2 and Figure 1(c), study an observed treatment-induced recanting witness \(D\) and a distinct latent baseline confounder \(U\) affecting \(A,D,M,Y\). Their Assumption 5 treats \(Z\) and \(W\) as proxies for \(U\), with \(W\mathbin{\perp\!\!\!\perp}\{A,D,M\}\mid(U,X)\) and \(Z\mathbin{\perp\!\!\!\perp}\{W,D,M,Y\}\mid(A,U,X)\); in the depicted orientation, \(Z\) is treatment-inducing toward \(A\) and \(W\) is outcome-inducing toward \(Y\). Their Assumption 6(i)--(iii) imposes three stage-specific completeness or injectivity conditions for \(U\) through \(Z\). Under Assumptions 1--6, their Theorem 2.1 and equations (4)--(7) use outcome bridges \(h_2,h_1,h_0\) to identify \(\psi_B=\mathbb E[Y(1,D(1),M(0,D(1)))]\) by \(\mathbb E[h_0(W,X)]\). Their discussion after Theorem 2.1 explicitly permits nonunique bridge solutions; the contrast here concerns completeness for the causal or latent lift, not bridge uniqueness. The present graph instead makes the treatment-induced recanting witness \(H\) itself latent, generates \(Z\) from \((A,H)\), and generates \(W\) from \(H\) in the outcome system. Its target is \(\psi(\mathcal S)=\mathbb E[Y_{0,H_0,M_{1,H_0},W_{H_0}}]\): Bai--Wu--Sun--Cui hold the observed recanting witness at its treatment-one realization and import the treatment-zero mediator, whereas this proposal holds the hidden witness and outcome regime at zero and imports the treatment-one mediator while retaining the same unit-level \(H_0\) and \(W_{H_0}\). Treatment-label reversal alone does not remove the graph, latent-witness, proxy-target, or assumption differences. Bai--Wu--Sun--Cui do not subsume Theorem 1 because their recanting witness is observed and distinct from \(U\), Assumption 6 requires three stage-specific completeness conditions, and Theorem 2.1 targets \(\psi_B\) on that graph; the exact positive WBC success member is rank deficient at both relevant three-state and binary-proxy maps and lies outside that completeness route. Theorem 1 does not subsume Bai--Wu--Sun--Cui because it treats only the fixed hidden-\(H\) WBC graph and scalar \(\psi(\mathcal S)\), relies on the calibrated common-shock subclass or an explicitly assumed latent-valid bridge intersection plus two non-testable latent transport-overlap restrictions, and does not cover their observed-\(D\), baseline-\(U\) graph, \(\psi_B\), general three-bridge path-specific construction, four identification strategies, or inference results. Duarte--Finkelstein--Knox--Mummolo--Shpitser (2024), Section 4, Algorithm 2, and Theorem 1 establish the generic response-function polynomial-program reduction and identify its extrema as sharp bounds for discrete causal problems; their Algorithm 3 supplies anytime \(\varepsilon\)-sharp optimization. This paper credits Duarte et al. for that generic reduction and optimization machinery. The new content claimed here is only the exact WBC bridge-supported causal-fiber equality, its maximal-class specialization, the causal bridge and rescue results, the explicit WBC success and failure certificates, and fixed-channel one-bit completeness restoration. This paper does not prove a new generic response-program reduction, terminating endpoint solver, branch-and-bound or sum-of-squares certificate theorem, endpoint attainment, certified approximation, or inference theorem. Conversely, Duarte et al. do not supply the WBC-specific bridge-supported class, below-completeness calibrated causal lift, proxy-measurable rescue, exact WBC success and failure artifacts, or fixed-channel one-bit completeness restoration. Zhang--Li--Miao--Tchetgen Tchetgen (2023), Propositions 2.3 and 2.5, supply the inherited null-annihilator equivalence and distinguish identification from root-\(n\) estimability. Shpitser--Wood-Doughty--Tchetgen Tchetgen (2023) Assumption 2.5 requires injectivity of \(v(U)\mapsto\mathbb E[v(U)\mid Z,a,C]\) in the ordinary proximal problem; Assumption 5.13 requires the corresponding sequential injectivity of \(v(U^\star)\mapsto\mathbb E[v(U^\star)\mid a,T,Z,\operatorname{do}(W)]\) at a recursive proximal-fixing step. Under sequential proxy independence, fixing ignorability, bridge existence, sequential completeness, and positivity, their Theorem 4 identifies the next inductive margin \(p(R,T\mid\operatorname{do}(W,a))\). Their Section 5.3 discussion applies Proximal ID only on graph-model subsets where the required Fredholm equations have solutions and leaves completeness of the algorithm as an identification characterization open. Proximal ID does not subsume the present theorem: the exact calibrated three-state and binary-proxy success SCM has nontrivial null directions in both relevant proxy maps, while the scalar shared-witness target is invariant by the calibrated bridge and loading-row-space condition, and Theorem 4 does not identify this same-unit hidden-recanting mean. The present theorem does not subsume Proximal ID: it treats one fixed WBC graph and one scalar cross-world target, assumes WBC-specific causal and latent transport restrictions, and supplies neither recursive graphical fixing nor general distributional margins. The present result, Bai--Wu--Sun--Cui, and Proximal ID are conditional in different senses; no one of them is used here as a complete identification characterization of the others' graph and target. Bennett et al. (2023), Kallus--Mao--Uehara (2021 manuscript; 2024 journal version), Santos (2011), and Severini--Tripathi (2012) concern strongly identified inverse-problem functionals but not this WBC cross-world causal lift. Miao--Geng--Tchetgen Tchetgen (2018) supplies finite proxy-rank logic. Ghassami--Zhang--Shpitser--Tchetgen Tchetgen (2026, version 4) supplies incomplete-proxy bounds, including hidden-mediator analogues, rather than this scalar identification result. Dukes--Shpitser--Tchetgen Tchetgen (2023) studies proximal mediation under a different sequential structure. The present contribution is the primitive calibrated causal lift, the exact WBC causal-fiber specialization, and full-SCM constructions, not the generic response-program reduction or the stacking of two bridge stages.

\section{Setup}

Target (estimand / structural parameter): \(\psi(\mathcal S)=\mathbb E_{\mathcal S}[Y_{0,H_0,M_{1,H_0},W_{H_0}}]\).
Primary functional / estimator / diagnostic: For \(P\) induced by \(\mathfrak M_{\mathrm{WBC}}^{+}\), the canonical full-law range is \(\Gamma(P)=\{T(\theta):(\theta,x)\in\mathcal F(P)\}\), with \(\underline\psi(P)=\inf\Gamma(P)\) and \(\overline\psi(P)=\sup\Gamma(P)\), and the causal target is point identified on its full-law fiber if and only if \(\underline\psi(P)=\overline\psi(P)\). These are exact population optimization definitions and characterizations; they do not assert endpoint attainment or a terminating computational procedure. As sufficient point-identification routes, the row-space condition gives \(\psi(\mathcal S)=\ell(P)^{\top}x=r(P)^{\top}b(P)\) when \(B(P)x=b(P)\) and \(B(P)^{\top}r(P)=\ell(P)\), while the proxy-measurable arm-zero condition gives \(\psi(\mathcal S)=\mathbb E_P[f(W)\mid A=0]\) whenever \(P(Y=f(W)\mid A=0)=1\), without the row-space condition.
\begin{itemize}
  \item \(\mathcal A\) --- \texttt{finite alphabet}; \(\{0,1\}\); \texttt{alphabet}
  \par\smallskip\noindent\emph{Definition.} \texttt{Treatment alphabet.}
  \item \(\mathcal H\) --- \texttt{finite alphabet}; \(\{0,1,2\}\); \texttt{alphabet}
  \par\smallskip\noindent\emph{Definition.} \texttt{Hidden-\allowbreak{}witness alphabet.}
  \item \(\mathcal X\) --- \texttt{finite alphabet}; \(\{0,1\}\); \texttt{alphabet}
  \par\smallskip\noindent\emph{Definition.} \texttt{Common observed post-\allowbreak{}treatment alphabet.}
  \item \(a\) --- \texttt{treatment index}; \(\mathcal A\); \texttt{index}
  \par\smallskip\noindent\emph{Definition.} \texttt{Generic treatment level.}
  \item \(h\) --- \texttt{hidden-\allowbreak{}state index}; \(\mathcal H\); \texttt{index}
  \par\smallskip\noindent\emph{Definition.} \texttt{Generic hidden state.}
  \item \(z\) --- \texttt{proxy-\allowbreak{}state index}; \(\mathcal X\); \texttt{index}
  \par\smallskip\noindent\emph{Definition.} \texttt{Generic treatment-\allowbreak{}proxy state.}
  \item \(w\) --- \texttt{proxy-\allowbreak{}state index}; \(\mathcal X\); \texttt{index}
  \par\smallskip\noindent\emph{Definition.} \texttt{Generic outcome-\allowbreak{}proxy state.}
  \item \(m\) --- \texttt{mediator-\allowbreak{}state index}; \(\mathcal X\); \texttt{index}
  \par\smallskip\noindent\emph{Definition.} \texttt{Generic mediator state.}
  \item \(A\) --- \texttt{observed variable}; \(\mathcal A\); \texttt{SCM node}
  \par\smallskip\noindent\emph{Definition.} \texttt{Binary treatment.}
  \item \(H\) --- \texttt{latent variable}; \(\mathcal H\); \texttt{SCM node}
  \par\smallskip\noindent\emph{Definition.} \texttt{Three-\allowbreak{}state hidden recanting witness.}
  \item \(Z\) --- \texttt{observed variable}; \(\mathcal X\); \texttt{SCM node}
  \par\smallskip\noindent\emph{Definition.} \texttt{Binary treatment-\allowbreak{}inducing proxy.}
  \item \(W\) --- \texttt{observed variable}; \(\mathcal X\); \texttt{SCM node}
  \par\smallskip\noindent\emph{Definition.} \texttt{Binary outcome-\allowbreak{}inducing proxy.}
  \item \(M\) --- \texttt{observed variable}; \(\mathcal X\); \texttt{SCM node}
  \par\smallskip\noindent\emph{Definition.} Binary subsequent mediator, denoted \(M_2\) by Wu--Bai--Cui.
  \item \(Y\) --- \texttt{observed variable}; \(\mathcal X\); \texttt{SCM node}
  \par\smallskip\noindent\emph{Definition.} \texttt{Binary outcome.}
  \item \(O\) --- \texttt{observed vector}; \(\mathcal A\times\mathcal X^4\); \texttt{observed data}
  \par\smallskip\noindent\emph{Definition.} \(O=(A,Z,W,M,Y)\).
  \item \(\mathcal G_{\mathrm{WBC}}\) --- \texttt{acyclic directed mixed graph}; finite graphs on \(\{A,H,Z,W,M,Y\}\); \texttt{causal graph}
  \par\smallskip\noindent\emph{Definition.} Directed edges are \(A\to H,Z,M,Y\), \(H\to Z,W,M,Y\), \(M\to Y\), and \(W\to Y\); districts are \(\{\{A\},\{H,Y\},\{Z\},\{W\},\{M\}\}\).
  \item \(\mathcal S\) --- \texttt{finite structural causal model}; SCMs compatible with \(\mathcal G_{\mathrm{WBC}}\); \texttt{causal law}
  \par\smallskip\noindent\emph{Definition.} \texttt{Generic WBC SCM with jointly defined potential outcomes.}
  \item \(P_{\mathcal S}\) --- \texttt{full-\allowbreak{}data probability law}; laws induced by \(\mathcal S\); \texttt{full-\allowbreak{}data law}
  \par\smallskip\noindent\emph{Definition.} Joint law of observed, latent, and counterfactual variables induced by \(\mathcal S\).
  \item \(H_a\) --- \texttt{potential outcome}; \(\mathcal H\); \texttt{counterfactual}
  \par\smallskip\noindent\emph{Definition.} Hidden witness under \(A=a\).
  \item \(Z_{a,h}\) --- \texttt{potential outcome}; \(\mathcal X\); \texttt{counterfactual}
  \par\smallskip\noindent\emph{Definition.} Treatment-inducing proxy under \(A=a,H=h\).
  \item \(W_h\) --- \texttt{potential outcome}; \(\mathcal X\); \texttt{counterfactual}
  \par\smallskip\noindent\emph{Definition.} Outcome-inducing proxy under \(H=h\).
  \item \(M_{a,h}\) --- \texttt{potential outcome}; \(\mathcal X\); \texttt{counterfactual}
  \par\smallskip\noindent\emph{Definition.} Mediator under \(A=a,H=h\).
  \item \(Y_{a,h,m,w}\) --- \texttt{potential outcome}; \(\mathcal X\); \texttt{counterfactual}
  \par\smallskip\noindent\emph{Definition.} Outcome under \(A=a,H=h,M=m,W=w\).
  \item \(\psi(\mathcal S)\) --- \texttt{causal mean}; \([0,1]\); \texttt{target estimand}
  \par\smallskip\noindent\emph{Definition.} \(\psi(\mathcal S)=\mathbb E_{\mathcal S}[Y_{0,H_0,M_{1,H_0},W_{H_0}}]\), with the two appearances of \(H_0\) denoting the same unit-level realization.
  \item \(P_{\mathcal S}^O\) --- \texttt{probability law}; the \(31\)-simplex; \texttt{observed-\allowbreak{}law map}
  \par\smallskip\noindent\emph{Definition.} Observed law of \(O\) induced by \(\mathcal S\).
  \item \(P\) --- \texttt{probability law}; the \(31\)-simplex; \texttt{observed law}
  \par\smallskip\noindent\emph{Definition.} Generic observed law of \(O\).
  \item \(h_0\) --- \texttt{outcome bridge}; \(\mathbb R^{\mathcal X\times\mathcal X}\); \texttt{bridge nuisance}
  \par\smallskip\noindent\emph{Definition.} First-stage bridge indexed by \(W,M\).
  \item \(h_1\) --- \texttt{outcome bridge}; \(\mathbb R^{\mathcal X}\); \texttt{bridge nuisance}
  \par\smallskip\noindent\emph{Definition.} Second-stage bridge indexed by \(W\).
  \item \(x\) --- \texttt{bridge vector}; \(\mathbb R^6\); \texttt{finite bridge coordinate}
  \par\smallskip\noindent\emph{Definition.} \(x=(h_0(0,0),h_0(1,0),h_0(0,1),h_0(1,1),h_1(0),h_1(1))^{\top}\).
  \item \(B(P)\) --- \texttt{observable matrix}; \(\mathbb R^{6\times6}\); \texttt{observable operator}
  \par\smallskip\noindent\emph{Definition.} Coupled bridge matrix in \(\mathrm{def{:}coupled\text{-}bridge\text{-}matrix}\).
  \item \(b(P)\) --- \texttt{observable vector}; \(\mathbb R^6\); \texttt{observable moment vector}
  \par\smallskip\noindent\emph{Definition.} Right side in \(\mathrm{def{:}coupled\text{-}bridge\text{-}matrix}\).
  \item \(\ell(P)\) --- \texttt{observable loading}; \(\mathbb R^6\); \texttt{target loading}
  \par\smallskip\noindent\emph{Definition.} \(\ell(P)=(0,0,0,0,P(W=0\mid A=0),P(W=1\mid A=0))^{\top}\).
  \item \(\mathcal H_{\mathrm{obs}}(P)\) --- \texttt{affine bridge fiber}; subsets of \(\mathbb R^6\); \texttt{observable solution set}
  \par\smallskip\noindent\emph{Definition.} \(\{x:B(P)x=b(P)\}\).
  \item \(\mathcal H_{\mathrm{lat}}(\mathcal S)\) --- \texttt{latent-\allowbreak{}valid bridge set}; subsets of \(\mathbb R^6\); \texttt{causal bridge set}
  \par\smallskip\noindent\emph{Definition.} Sequential latent bridge pairs in \(\mathrm{def{:}latent\text{-}bridge\text{-}set}\).
  \item \(\mathfrak M_{\mathrm{WBC}}^{+}\) --- \texttt{SCM class}; subsets of SCMs on \(\mathcal G_{\mathrm{WBC}}\); \texttt{model class}
  \par\smallskip\noindent\emph{Definition.} Bridge-supported below-completeness class in \(\mathrm{def{:}wbc\text{-}class}\).
  \item \(\mathsf{RI}(P)\) --- \texttt{observable predicate}; \(\{\mathrm{true},\mathrm{false}\}\); \texttt{identification criterion}
  \par\smallskip\noindent\emph{Definition.} Observed solvability plus target loading in the row space of \(B(P)\).
  \item \(r(P)\) --- \texttt{row-\allowbreak{}space representer}; \(\mathbb R^6\); \texttt{identifying weight}
  \par\smallskip\noindent\emph{Definition.} Any solution of \(B(P)^{\top}r(P)=\ell(P)\).
  \item \(\mathcal S^{+}\) --- \texttt{finite WBC SCM}; SCMs on \(\mathcal G_{\mathrm{WBC}}\); \texttt{success witness}
  \par\smallskip\noindent\emph{Definition.} \texttt{Exact success artifact.}
  \item \(P^{+}\) --- \texttt{probability law}; the \(31\)-simplex; \texttt{success law}
  \par\smallskip\noindent\emph{Definition.} Observed law of \(\mathcal S^+\).
  \item \(\mathcal S^{\oplus}\) --- \texttt{finite WBC SCM}; SCMs on \(\mathcal G_{\mathrm{WBC}}\); \texttt{counterpair member}
  \par\smallskip\noindent\emph{Definition.} \texttt{First failure SCM.}
  \item \(\mathcal S^{\ominus}\) --- \texttt{finite WBC SCM}; SCMs on \(\mathcal G_{\mathrm{WBC}}\); \texttt{counterpair member}
  \par\smallskip\noindent\emph{Definition.} \texttt{Second failure SCM.}
  \item \(P^{\dagger}\) --- \texttt{probability law}; the \(31\)-simplex; \texttt{failure law}
  \par\smallskip\noindent\emph{Definition.} \texttt{Common observed law of the counterpair.}
  \item \(V\) --- \texttt{observed variable}; \(\mathcal X\); \texttt{augmented proxy}
  \par\smallskip\noindent\emph{Definition.} \texttt{Additional binary treatment-\allowbreak{}inducing proxy.}
  \item \(Z^{\star}\) --- \texttt{joint proxy}; \(\mathcal X^2\); \texttt{augmented treatment proxy}
  \par\smallskip\noindent\emph{Definition.} \(Z^{\star}=(Z,V)\).
  \item \(R_{Z^{\star}\mid H}\) --- \texttt{proxy-\allowbreak{}channel matrix}; \([0,1]^{4\times3}\); \texttt{augmentation operator}
  \par\smallskip\noindent\emph{Definition.} Channel of \(Z^{\star}\) given \(H\).
  \item \(\mathsf{Tilt}_{\mathrm{WBC}}\) --- \texttt{construction handle}; \texttt{latent perturbation schemes}; \texttt{open-\allowbreak{}question handle}
  \par\smallskip\noindent\emph{Definition.} \texttt{Interior or complement-\allowbreak{}of-\allowbreak{}rescue converse handle.}
  \item \(q\) --- \texttt{calibrated severity vector}; \([0,1]^3\); \texttt{sufficient-\allowbreak{}submodel parameter}
  \par\smallskip\noindent\emph{Definition.} The vector with coordinate \(q_h=P(W_h=1)\).
  \item \(\alpha\) --- \texttt{real parameter}; \(\mathbb R\); \texttt{sufficient-\allowbreak{}submodel parameter}
  \par\smallskip\noindent\emph{Definition.} \texttt{Outcome intercept in the calibrated sufficient submodel.}
  \item \(\tau\) --- \texttt{real parameter}; \(\mathbb R\); \texttt{sufficient-\allowbreak{}submodel parameter}
  \par\smallskip\noindent\emph{Definition.} \texttt{Additive mediator coefficient in the calibrated sufficient submodel.}
  \item \(\beta\) --- \texttt{real parameter}; \(\mathbb R\); \texttt{sufficient-\allowbreak{}submodel parameter}
  \par\smallskip\noindent\emph{Definition.} \texttt{Additive measured-\allowbreak{}severity coefficient in the calibrated sufficient submodel.}
  \item \(c\) --- \texttt{real parameter}; \(\mathbb R\); \texttt{sufficient-\allowbreak{}submodel parameter}
  \par\smallskip\noindent\emph{Definition.} \texttt{Mediator-\allowbreak{}risk intercept in the calibrated sufficient submodel.}
  \item \(d\) --- \texttt{real parameter}; \(\mathbb R\); \texttt{sufficient-\allowbreak{}submodel parameter}
  \par\smallskip\noindent\emph{Definition.} \texttt{Mediator-\allowbreak{}risk slope in the calibrated sufficient submodel.}
  \item \(\Theta_{\mathrm{cal}}\) --- \texttt{parameter set}; subsets of \([0,1]^3\times\mathbb R^5\); \texttt{sufficient-\allowbreak{}submodel parameter space}
  \par\smallskip\noindent\emph{Definition.} Tuples \((q,\alpha,\tau,\beta,c,d)\) such that \(0\le c+dq_h\le1\) for every \(h\) and \(0\le\alpha+\tau m+\beta w\le1\) for every \((m,w)\in\mathcal X^2\).
  \item \(U_W\) --- \texttt{exogenous variable}; \([0,1]\); \texttt{sufficient-\allowbreak{}submodel shock}
  \par\smallskip\noindent\emph{Definition.} \texttt{Uniform shock for the calibrated outcome-\allowbreak{}proxy channel.}
  \item \(U_M\) --- \texttt{exogenous variable}; \([0,1]\); \texttt{sufficient-\allowbreak{}submodel shock}
  \par\smallskip\noindent\emph{Definition.} \texttt{Uniform shock for the calibrated treatment-\allowbreak{}one mediator.}
  \item \(U_Y\) --- \texttt{exogenous variable}; \([0,1]\); \texttt{sufficient-\allowbreak{}submodel shock}
  \par\smallskip\noindent\emph{Definition.} \texttt{Uniform shock for the calibrated treatment-\allowbreak{}zero outcome.}
  \item \(\mathcal S^{\mathrm{cal}}\) --- \texttt{parameterized finite WBC SCM}; SCMs on \(\mathcal G_{\mathrm{WBC}}\); \texttt{sufficient-\allowbreak{}submodel witness}
  \par\smallskip\noindent\emph{Definition.} Generic member of the calibrated sufficient family constructed in \(\mathrm{def{:}calibrated\text{-}bridge\text{-}family\).
  \item \(\mathfrak M_{\mathrm{cal}}\) --- \texttt{SCM subclass}; subsets of SCMs on \(\mathcal G_{\mathrm{WBC}}\); \texttt{principal below-\allowbreak{}completeness model class}
  \par\smallskip\noindent\emph{Definition.} Calibrated primitive WBC subclass in \(\mathrm{def{:}calibrated\text{-}wbc\text{-}subclass}\).
\end{itemize}

\section{Assumptions}

\begin{assumption}[ass:treatment-consistency]\label{ass:treatment-consistency}
\(H=H_A\) almost surely.
\par\smallskip\noindent\emph{Standard} (WBC hidden-witness consistency, \cite{WuBaiCui2026}).
\end{assumption}

\begin{assumption}[ass:z-consistency]\label{ass:z-consistency}
\(Z=Z_{A,H}\) almost surely.
\par\smallskip\noindent\emph{Standard} (Proxy consistency, \cite{WuBaiCui2026}).
\end{assumption}

\begin{assumption}[ass:w-consistency]\label{ass:w-consistency}
\(W=W_H\) almost surely.
\par\smallskip\noindent\emph{Standard} (Proxy consistency, \cite{WuBaiCui2026}).
\end{assumption}

\begin{assumption}[ass:m-consistency]\label{ass:m-consistency}
\(M=M_{A,H}\) almost surely.
\par\smallskip\noindent\emph{Standard} (Mediator consistency, \cite{WuBaiCui2026}).
\end{assumption}

\begin{assumption}[ass:y-consistency]\label{ass:y-consistency}
\(Y=Y_{A,H,M,W}\) almost surely.
\par\smallskip\noindent\emph{Standard} (Outcome consistency, \cite{WuBaiCui2026}).
\end{assumption}

\begin{assumption}[ass:treatment-positivity]\label{ass:treatment-positivity}
\(0<P_{\mathcal S}^O(A=a)<1\) for every \(a\in\mathcal A\).
\par\smallskip\noindent\emph{Standard} (WBC treatment positivity, finite no-covariate version, \cite{WuBaiCui2026}).
\end{assumption}

\begin{assumption}[ass:a0-bridge-denominators]\label{ass:a0-bridge-denominators}
\(P_{\mathcal S}^O(A=0,Z=z,M=m)>0\) for every \((z,m)\in\mathcal X^2\).
\par\smallskip\noindent\emph{Standard} (Finite conditional-moment denominator support, \cite{ZhangLiMiaoTchetgen2023}).
\end{assumption}

\begin{assumption}[ass:a1-bridge-denominators]\label{ass:a1-bridge-denominators}
\(P_{\mathcal S}^O(A=1,Z=z)>0\) for every \(z\in\mathcal X\).
\par\smallskip\noindent\emph{Standard} (Finite conditional-moment denominator support, \cite{ZhangLiMiaoTchetgen2023}).
\end{assumption}

\begin{assumption}[ass:sequential-treatment]\label{ass:sequential-treatment}
\(\{Y_{a,H_a,m,W_{H_a}},H_a\}\mathbin{\perp\!\!\!\perp}A\) for all \(a,m\).
\par\smallskip\noindent\emph{Standard} (WBC sequential ignorability, treatment component, \cite{WuBaiCui2026}).
\end{assumption}

\begin{assumption}[ass:sequential-outcome]\label{ass:sequential-outcome}
\(Y_{A,H,m,W_H}\mathbin{\perp\!\!\!\perp}M\mid(H,A)\) for every \(m\in\mathcal X\).
\par\smallskip\noindent\emph{Standard} (WBC sequential ignorability, outcome component, \cite{WuBaiCui2026}).
\end{assumption}

\begin{assumption}[ass:sequential-mediator]\label{ass:sequential-mediator}
\(M_{a,h}\mathbin{\perp\!\!\!\perp}(H,A)\) for every \((a,h)\in\mathcal A\times\mathcal H\).
\par\smallskip\noindent\emph{Standard} (WBC sequential ignorability, mediator component, \cite{WuBaiCui2026}).
\end{assumption}

\begin{assumption}[ass:cross-world]\label{ass:cross-world}
\(\{Y_{a,H_a,m,W_{H_a}},H_a\}\mathbin{\perp\!\!\!\perp}M_{a',h}\) for every \((a,a',h,m)\).
\par\smallskip\noindent\emph{Standard} (WBC cross-world independence, \cite{WuBaiCui2026}).
\end{assumption}

\begin{assumption}[ass:z-proxy-exclusion]\label{ass:z-proxy-exclusion}
\(Z\mathbin{\perp\!\!\!\perp}(Y,M)\mid(H,A)\).
\par\smallskip\noindent\emph{Standard} (WBC treatment-inducing proxy exclusion, \cite{WuBaiCui2026}).
\end{assumption}

\begin{assumption}[ass:w-proxy-exclusion]\label{ass:w-proxy-exclusion}
\(W\mathbin{\perp\!\!\!\perp}(A,Z,M)\mid H\).
\par\smallskip\noindent\emph{Standard} (WBC outcome-inducing proxy exclusion, \cite{WuBaiCui2026}).
\end{assumption}

\begin{assumption}[ass:latent-a0-transport-support]\label{ass:latent-a0-transport-support}
For every \((h,m)\in\mathcal H\times\mathcal X\) with \(P_{\mathcal S}(H_0=h)>0\) and \(P_{\mathcal S}(M_{1,h}=m)>0\), \(P_{\mathcal S}(H=h,M=m,A=0)>0\).
\par\smallskip\noindent\emph{Novel.} This proposal-specific cross-world support restriction is defensible when, for every target-relevant latent state \(h\), no mediator category possible under treatment one is structurally impossible under treatment zero at that state. A sufficient primitive support scenario is \(\operatorname{supp}(M_{1,h})\subseteq\operatorname{supp}(M_{0,h})\), together with maintained consistency, treatment and mediator independence, and positive arm sampling. In the intended school-track application, every achievement category possible under the academic track for a relevant latent non-cognitive type is represented among otherwise comparable general-track units. This latent restriction is not testable from \(O\) and must be defended from design or subject-matter knowledge.
\end{assumption}

\begin{assumption}[ass:latent-a1-transport-support]\label{ass:latent-a1-transport-support}
For every \(h\in\mathcal H\) with \(P_{\mathcal S}(H_0=h)>0\), \(P_{\mathcal S}(H=h,A=1)>0\).
\par\smallskip\noindent\emph{Novel.} This proposal-specific cross-world hidden-state overlap restriction is defensible when every latent type possible under treatment zero remains possible under treatment one. A sufficient primitive support scenario is \(\operatorname{supp}(H_0)\subseteq\operatorname{supp}(H_1)\), together with maintained consistency, treatment independence, and positive arm sampling. In the intended school-track application, it rules out deterministic exclusion of any target-relevant latent non-cognitive type from the academic-track arm. This latent restriction is not testable from \(O\) and must be defended from design or subject-matter knowledge.
\end{assumption}

\begin{assumption}[ass:latent-valid-bridge-existence]\label{ass:latent-valid-bridge-existence}
\(\mathcal H_{\mathrm{lat}}(\mathcal S)\cap\mathcal H_{\mathrm{obs}}(P_{\mathcal S}^O)\neq\varnothing\).
\par\smallskip\noindent\emph{Novel.} Latent sequential bridge realizability is credible when calibrated severity drives an affine treatment-one mediator risk and treatment-zero outcomes have additive mediator and measured-severity effects. The calibrated-subclass theorem places the explicit sequential bridge pair in the observed affine solution set, and both exact rational artifacts certify the broader restriction while three-state completeness fails.
\end{assumption}

\section{Definitions}

\begin{definition}[def:coupled-bridge-matrix]\label{def:coupled-bridge-matrix}
\par\smallskip\noindent\emph{Defined object.} \texttt{Observable coupled WBC bridge system}
\par\smallskip\noindent\emph{Construction.} When \(P(A=0,Z=z,M=m)>0\) for every \((z,m)\in\mathcal X^2\) and \(P(A=1,Z=z)>0\) for every \(z\in\mathcal X\), order the first four rows by \((z,m)=(0,0),(1,0),(0,1),(1,1)\) and the last two by \(z=0,1\). Set \(B(P)_{(0,z,m),(w',m')}=\mathbf 1\{m=m'\}P(W=w'\mid Z=z,M=m,A=0)\), \(B(P)_{(0,z,m),h_1(w')}=0\), \(b(P)_{(0,z,m)}=\mathbb E_P[Y\mid Z=z,M=m,A=0]\), \(B(P)_{(1,z),(w',m')}=P(W=w',M=m'\mid Z=z,A=1)\), \(B(P)_{(1,z),h_1(w')}=-P(W=w'\mid Z=z,A=1)\), and \(b(P)_{(1,z)}=0\). Every entry is observed, and no positivity of any value of \(W\) or \(Y\) is imposed.
\par\smallskip\noindent\emph{Inputs.} \(P\).
\end{definition}

\begin{definition}[def:observed-bridge-fiber]\label{def:observed-bridge-fiber}
\par\smallskip\noindent\emph{Defined object.} \texttt{Observed affine bridge fiber}
\par\smallskip\noindent\emph{Construction.} \(\mathcal H_{\mathrm{obs}}(P)=\{x\in\mathbb R^6:B(P)x=b(P)\}\).
\par\smallskip\noindent\emph{Inputs.} \(P\).
\end{definition}

\begin{definition}[def:latent-bridge-set]\label{def:latent-bridge-set}
\par\smallskip\noindent\emph{Defined object.} \texttt{Latent-\allowbreak{}valid sequential outcome bridges}
\par\smallskip\noindent\emph{Construction.} \(\mathcal H_{\mathrm{lat}}(\mathcal S)\) is the set of vectors induced by pairs \((h_0,h_1)\) satisfying \(\mathbb E_{P_{\mathcal S}}[Y\mid H=h,M=m,A=0]=\sum_w h_0(w,m)P_{\mathcal S}(W=w\mid H=h,M=m,A=0)\) for every \((h,m)\) with \(P_{\mathcal S}(H_0=h)>0\) and \(P_{\mathcal S}(M_{1,h}=m)>0\), and satisfying \(\sum_{w,m}h_0(w,m)P_{\mathcal S}(W=w,M=m\mid H=h,A=1)=\sum_w h_1(w)P_{\mathcal S}(W=w\mid H=h,A=1)\) for every \(h\) with \(P_{\mathcal S}(H_0=h)>0\). The two proposal-specific target-relevant latent transport-overlap restrictions make each displayed conditional law well defined; they are additional restrictions rather than standard WBC positivity premises.
\par\smallskip\noindent\emph{Inputs.} \(\mathcal S\).
\end{definition}

\begin{definition}[def:calibrated-bridge-family]\label{def:calibrated-bridge-family}
\par\smallskip\noindent\emph{Defined object.} \texttt{Calibrated binary-\allowbreak{}measurement sufficient family}
\par\smallskip\noindent\emph{Construction.} Choose \((q,\alpha,\tau,\beta,c,d)\in\Theta_{\mathrm{cal}}\). Let \(U_W,U_M,U_Y\) be mutually independent \(\operatorname{Unif}(0,1)\) variables, independent of the remaining exogenous variables, and retain finite WBC structural kernels satisfying the maintained consistency, sequential, cross-world, and exclusion conditions, together with the two proposal-specific target-relevant latent transport-overlap restrictions, with \(P_{\mathcal S^{\mathrm{cal}}}(H_0=h)>0\) for all three \(h\in\mathcal H\). Set \(W_h=\mathbf 1\{U_W\le q_h\}\), \(M_{1,h}=\mathbf 1\{U_M\le c+dq_h\}\), and \(Y_{0,h,m,w}=\mathbf 1\{U_Y\le\alpha+\tau m+\beta w\}\); this parameterized construction is \(\mathcal S^{\mathrm{cal}}\). Associate \(h_0(w,m)=\alpha+\tau m+\beta w\) and \(h_1(w)=\alpha+\tau c+(\beta+\tau d)w\). The construction does not derive either latent transport-overlap restriction and does not specify enough of \(H_1\) or \(M_{0,h}\) to do so.
\par\smallskip\noindent\emph{Inputs.} \(q\); \(\alpha\); \(\tau\); \(\beta\); \(c\); \(d\).
\end{definition}

\begin{definition}[def:calibrated-wbc-subclass]\label{def:calibrated-wbc-subclass}
\par\smallskip\noindent\emph{Defined object.} \texttt{Calibrated primitive WBC subclass}
\par\smallskip\noindent\emph{Construction.} \(\mathfrak M_{\mathrm{cal}}=\{\mathcal S^{\mathrm{cal}}:\mathcal S^{\mathrm{cal}}\text{ is generated by }\mathrm{def{:}calibrated\text{-}bridge\text{-}family\text{ and satisfies all listed member properties}\}\). The two target-relevant latent transport-overlap restrictions are retained member properties and are not consequences of the common-shock equations.
\par\smallskip\noindent\emph{Member properties.} \texttt{ass:\allowbreak{}treatment-\allowbreak{}consistency}; \texttt{ass:\allowbreak{}z-\allowbreak{}consistency}; \texttt{ass:\allowbreak{}w-\allowbreak{}consistency}; \texttt{ass:\allowbreak{}m-\allowbreak{}consistency}; \texttt{ass:\allowbreak{}y-\allowbreak{}consistency}; \texttt{ass:\allowbreak{}a0-\allowbreak{}bridge-\allowbreak{}denominators}; \texttt{ass:\allowbreak{}a1-\allowbreak{}bridge-\allowbreak{}denominators}; \texttt{ass:\allowbreak{}sequential-\allowbreak{}treatment}; \texttt{ass:\allowbreak{}sequential-\allowbreak{}outcome}; \texttt{ass:\allowbreak{}sequential-\allowbreak{}mediator}; \texttt{ass:\allowbreak{}cross-\allowbreak{}world}; \texttt{ass:\allowbreak{}z-\allowbreak{}proxy-\allowbreak{}exclusion}; \texttt{ass:\allowbreak{}w-\allowbreak{}proxy-\allowbreak{}exclusion}; \texttt{ass:\allowbreak{}latent-\allowbreak{}a0-\allowbreak{}transport-\allowbreak{}support}; \texttt{ass:\allowbreak{}latent-\allowbreak{}a1-\allowbreak{}transport-\allowbreak{}support}.
\end{definition}

\begin{definition}[def:row-space-criterion]\label{def:row-space-criterion}
\par\smallskip\noindent\emph{Defined object.} \texttt{Observable row-\allowbreak{}space criterion}
\par\smallskip\noindent\emph{Construction.} \(\mathsf{RI}(P)\Longleftrightarrow[b(P)\in\operatorname{col}(B(P))\text{ and }\ell(P)\in\operatorname{row}(B(P))]\).
\par\smallskip\noindent\emph{Inputs.} \(P\).
\end{definition}

\begin{definition}[def:wbc-class]\label{def:wbc-class}
\par\smallskip\noindent\emph{Defined object.} \texttt{Conditional bridge-\allowbreak{}supported WBC envelope}
\par\smallskip\noindent\emph{Construction.} \(\mathfrak M_{\mathrm{WBC}}^{+}=\{\mathcal S:\mathcal S\text{ is compatible with }\mathcal G_{\mathrm{WBC}}\text{ and satisfies all listed member properties}\}\).
\par\smallskip\noindent\emph{Member properties.} \texttt{ass:\allowbreak{}treatment-\allowbreak{}consistency}; \texttt{ass:\allowbreak{}z-\allowbreak{}consistency}; \texttt{ass:\allowbreak{}w-\allowbreak{}consistency}; \texttt{ass:\allowbreak{}m-\allowbreak{}consistency}; \texttt{ass:\allowbreak{}y-\allowbreak{}consistency}; \texttt{ass:\allowbreak{}a0-\allowbreak{}bridge-\allowbreak{}denominators}; \texttt{ass:\allowbreak{}a1-\allowbreak{}bridge-\allowbreak{}denominators}; \texttt{ass:\allowbreak{}sequential-\allowbreak{}treatment}; \texttt{ass:\allowbreak{}sequential-\allowbreak{}outcome}; \texttt{ass:\allowbreak{}sequential-\allowbreak{}mediator}; \texttt{ass:\allowbreak{}cross-\allowbreak{}world}; \texttt{ass:\allowbreak{}z-\allowbreak{}proxy-\allowbreak{}exclusion}; \texttt{ass:\allowbreak{}w-\allowbreak{}proxy-\allowbreak{}exclusion}; \texttt{ass:\allowbreak{}latent-\allowbreak{}a0-\allowbreak{}transport-\allowbreak{}support}; \texttt{ass:\allowbreak{}latent-\allowbreak{}a1-\allowbreak{}transport-\allowbreak{}support}; \texttt{ass:\allowbreak{}latent-\allowbreak{}valid-\allowbreak{}bridge-\allowbreak{}existence}.
\end{definition}

\begin{definition}[def:success-artifact]\label{def:success-artifact}
\par\smallskip\noindent\emph{Defined object.} \texttt{Exact rational rank-\allowbreak{}deficient success SCM}
\par\smallskip\noindent\emph{Construction.} Let independent \(U_A,U_H,U_Z,U_W,U_M,U_Y\sim\operatorname{Unif}(0,1)\) generate \(A=\mathbf 1\{U_A\le1/2\}\), \(H_a=0,1,2\) on intervals of lengths \(1/3,1/3,1/3\), \(Z_{a,h}=\mathbf 1\{U_Z\le q_h^Z\}\) with \(q^Z=(1/4,1/2,3/4)\), \(W_h=\mathbf 1\{U_W\le q_h^W\}\) with \(q^W=(2/3,1/6,2/3)\), \(M_{a,h}=\mathbf 1\{U_M\le1/2\}\), and \(Y_{a,h,m,w}=\mathbf 1\{U_Y\le1/2\}\). This is \(\mathcal S^+\), and \(P^+(a,z,w,m,y)=1/32\). Moreover, \(B(P^+)=\begin{pmatrix}1/2&1/2&0&0&0&0\\1/2&1/2&0&0&0&0\\0&0&1/2&1/2&0&0\\0&0&1/2&1/2&0&0\\1/4&1/4&1/4&1/4&-1/2&-1/2\\1/4&1/4&1/4&1/4&-1/2&-1/2\end{pmatrix}\), \(b(P^+)=(1/2,1/2,1/2,1/2,0,0)^{\top}\), and \(\ell(P^+)=(0,0,0,0,1/2,1/2)^{\top}\). Its rank is \(3\), a null basis is \(e_1-e_2,e_3-e_4,e_5-e_6\), and \(r(P^+)=(1/2,0,1/2,0,-1,0)^{\top}\). The distinct observed solutions \(x^{(1)}=(1,0,1,0,1,0)^{\top}\) and \(x^{(2)}=(0,1,0,1,0,1)^{\top}\) both have loading \(1/2\); \(x^{\mathrm{lat}}=(1/2,1/2,1/2,1/2,1/2,1/2)^{\top}\) is latent-valid and \(\psi(\mathcal S^+)=1/2\). The matrices \(P(H\mid Z)\) and \(P(H\mid W)\) have rank \(2\), with null vectors \((1,-2,1)^{\top}\) and \((1,0,-1)^{\top}\), respectively.
\end{definition}

\begin{definition}[def:failure-artifact]\label{def:failure-artifact}
\par\smallskip\noindent\emph{Defined object.} \texttt{Exact rational observationally equivalent WBC counterpair}
\par\smallskip\noindent\emph{Construction.} Use independent uniforms \(U_A,U_H,U_Z,U_W,U_M,U_Y\). In both SCMs set \(P(A=1)=1/2\). Using the same shared-\(U_H\) response construction in both \(\mathcal S^{\oplus}\) and \(\mathcal S^{\ominus}\), set \(H_0=0\) for \(U_H\in[0,1/2)\), \(H_0=1\) for \(U_H\in[1/2,5/6)\), and \(H_0=2\) for \(U_H\in[5/6,1]\); set \(H_1=0\) for \(U_H\in[0,1/3)\), \(H_1=1\) for \(U_H\in[1/3,2/3)\), and \(H_1=2\) for \(U_H\in[2/3,1]\). Endpoint conventions at the null boundaries are immaterial. These common response maps induce \(P(H_0=h)=(1/2,1/3,1/6)_h\) and \(P(H_1=h)=1/3\). Set \(Z_{a,h}=\mathbf 1\{U_Z\le1/2\}\), \(W_h=\mathbf 1\{U_W\le q_h\}\) with \(q=(1/4,1/2,3/4)\), and \(M_{0,h}=\mathbf 1\{U_M\le1/2\}\). For \(\mathcal S^{\oplus}\), let \(P(M_{1,h}=1)=(5/8,1/4,5/8)_h\); for \(\mathcal S^{\ominus}\), use \((3/8,3/4,3/8)_h\). In both, let \(P(Y_{0,h,0,w}=1)=1/2\), \(P(Y_{0,h,1,w}=1)=(0,1/2,1)_h\), and \(P(Y_{1,h,m,w}=1)=1/2\), using \(U_Y\). Their common law is \(P^{\dagger}(a,z,w,m,y)=\frac12\sum_{h=0}^2\pi_a(h)\frac12 q_h^w(1-q_h)^{1-w}s_{a,h}^m(1-s_{a,h})^{1-m}g_{a,h,m}^y(1-g_{a,h,m})^{1-y}\); the two \(s_{1,h}\) choices give the same sum because \((1,-2,1)^{\top}\) annihilates both \(W\)-channel columns. Every cell is positive. Here \(B(P^{\dagger})=\begin{pmatrix}7/12&5/12&0&0&0&0\\7/12&5/12&0&0&0&0\\0&0&7/12&5/12&0&0\\0&0&7/12&5/12&0&0\\1/4&1/4&1/4&1/4&-1/2&-1/2\\1/4&1/4&1/4&1/4&-1/2&-1/2\end{pmatrix}\), \(b(P^{\dagger})=(1/2,1/2,1/3,1/3,0,0)^{\top}\), and \(\ell(P^{\dagger})=(0,0,0,0,7/12,5/12)^{\top}\). The rank is \(3\), but \(\ell(P^{\dagger})^{\top}(0,0,0,0,1,-1)^{\top}=1/6\). Latent-valid vectors are \(x^{\oplus}=(1/2,1/2,-1/2,3/2,-1/8,9/8)^{\top}\) and \(x^{\ominus}=(1/2,1/2,-1/2,3/2,1/8,7/8)^{\top}\), yielding \(\psi(\mathcal S^{\oplus})=19/48\) and \(\psi(\mathcal S^{\ominus})=7/16\).
\end{definition}

\begin{definition}[def:binary-proxy-augmentation]\label{def:binary-proxy-augmentation}
\par\smallskip\noindent\emph{Defined object.} \texttt{One-\allowbreak{}bit treatment-\allowbreak{}proxy augmentation}
\par\smallskip\noindent\emph{Construction.} Order the rows by \((Z,V)=(0,0),(0,1),(1,0),(1,1)\) and the columns by \(H=0,1,2\). For the channel \(P(Z=1\mid H=h,A=a)=(1/4,1/2,3/4)_h\), add \(V\) with \(P(V=1\mid H=h,A=a)=(1/4,3/4,1/2)_h\), impose \(V\mathbin{\perp\!\!\!\perp}(Z,W,M,Y)\mid(H,A)\), and set \(Z^{\star}=(Z,V)\). Then \(R_{Z^{\star}\mid H}=\begin{pmatrix}9/16&1/8&1/8\\3/16&3/8&1/8\\3/16&1/8&3/8\\1/16&3/8&3/8\end{pmatrix}\). In particular, \(P(Z=1,V=0\mid H=0,A=a)=3/16\) and \(P(Z=1,V=1\mid H=0,A=a)=1/16\). The top-three-row minor satisfies \(\det\!\begin{pmatrix}9&2&2\\3&6&2\\3&2&6\end{pmatrix}/16^3=240/4096=15/256\neq0\).
\par\smallskip\noindent\emph{Inputs.} \(Z\); \(V\).
\end{definition}

\begin{definition}[def:converse-handle]\label{def:converse-handle}
\par\smallskip\noindent\emph{Defined object.} \texttt{Interior or complement-\allowbreak{}of-\allowbreak{}rescue latent null-\allowbreak{}tilt handle}
\par\smallskip\noindent\emph{Construction.} Given \(\mathcal S\in\mathfrak M_{\mathrm{WBC}}^{+}\) with \(P=P_{\mathcal S}^O\) and \(\ell(P)\notin\operatorname{row}(B(P))\), use \(\mathsf{Tilt}_{\mathrm{WBC}}\) only on either the strictly positive observed-law interior, where \(P(O=o)>0\) for every \(o\in\mathcal A\times\mathcal X^4\), or, more generally, the complement of the proxy-measurable treatment-zero rescue class, where \(P(Y=f(W)\mid A=0)<1\) for every map \(f:\mathcal X\to\mathcal X\). Seek a signed perturbation of the three-state \(M_{1,h}\)-kernel in the observational null space of the joint \((Z,W,M,Y)\) mixture map and a corresponding latent-valid sequential bridge displacement with nonzero \(\ell(P)\)-loading, subject to both perturbed objects being members of \(\mathfrak M_{\mathrm{WBC}}^{+}\), inducing the same \(P\), and having different values of \(\psi\). This is a search handle for an interior or complement-of-rescue converse; it asserts neither existence of such a perturbation nor maximality of the rescue class.
\par\smallskip\noindent\emph{Inputs.} \(P\); \(\mathcal S\).
\end{definition}

\section{Main results}

\begin{lemma}[lem:intersection-implies-solvability]\label{lem:intersection-implies-solvability}
For every \(\mathcal S\), if \(\mathcal H_{\mathrm{lat}}(\mathcal S)\cap\mathcal H_{\mathrm{obs}}(P_{\mathcal S}^O)\neq\varnothing\), then \(\mathcal H_{\mathrm{obs}}(P_{\mathcal S}^O)\neq\varnothing\), and therefore \(b(P_{\mathcal S}^O)\in\operatorname{col}(B(P_{\mathcal S}^O))\).
\end{lemma}
\begin{proof}
Let \(x_\star\in\mathcal H_{\mathrm{lat}}(\mathcal S)\cap
\mathcal H_{\mathrm{obs}}(P_{\mathcal S}^O)\), whose existence is the premise in
\(\mathrm{ass{:}latent\text{-}valid\text{-}bridge\text{-}existence}\).  By
\(\mathrm{def{:}observed\text{-}bridge\text{-}fiber}\),
\[
 B(P_{\mathcal S}^O)x_\star=b(P_{\mathcal S}^O).
\]
Thus \(x_\star\in\mathcal H_{\mathrm{obs}}(P_{\mathcal S}^O)\), so that set is
nonempty, and the displayed equality says exactly that
\(b(P_{\mathcal S}^O)\in\operatorname{col}(B(P_{\mathcal S}^O))\).
\end{proof}
\par\smallskip\noindent \textit{Justification.} The implication removes a redundant global observed-solvability premise from the conditional envelope. \textit{Gap.} This is elementary affine-set logic and carries no novelty claim. \textit{Consumer.} The conditional-envelope theorem needs only the loading row-space condition because solvability follows from latent-valid intersection.

\begin{lemma}[lem:null-annihilator]\label{lem:null-annihilator}
For every \(P\) with \(b(P)\in\operatorname{col}(B(P))\), the map \(x\mapsto\ell(P)^{\top}x\) is constant on \(\mathcal H_{\mathrm{obs}}(P)\) if and only if \(\ell(P)\in\operatorname{row}(B(P))\), equivalently if and only if \(\ell(P)^{\top}v=0\) for every \(v\in\operatorname{null}(B(P))\). If \(B(P)^{\top}r(P)=\ell(P)\), the common value is \(r(P)^{\top}b(P)\). This is the finite coupled specialization of Zhang--Li--Miao--Tchetgen Tchetgen, Propositions 2.3 and 2.5, and carries no novelty claim.
\end{lemma}
\begin{proof}
Fix \(P\) with \(b(P)\in\operatorname{col}(B(P))\), and choose
\(x_0\) with \(B(P)x_0=b(P)\).  The definition of the observed bridge fiber gives
\[
 \mathcal H_{\mathrm{obs}}(P)=x_0+\operatorname{null}(B(P)).
\]
Consequently, \(x\mapsto\ell(P)^\top x\) is constant on this affine set if and
only if
\[
 \ell(P)^\top v=0\qquad\text{for every }v\in\operatorname{null}(B(P)).
\]
For completeness, the finite-dimensional identity used here is derived directly.
Every row of \(B(P)\) annihilates \(\operatorname{null}(B(P))\), so
\(\operatorname{row}(B(P))\subseteq\operatorname{null}(B(P))^\perp\).  Rank--nullity
gives
\[
 \dim\operatorname{null}(B(P))^\perp
 =6-\dim\operatorname{null}(B(P))
 =\operatorname{rank}(B(P))
 =\dim\operatorname{row}(B(P)),
\]
and hence equality of the two subspaces.  This proves both equivalences.  Finally,
if \(B(P)^\top r(P)=\ell(P)\), then for every
\(x\in\mathcal H_{\mathrm{obs}}(P)\),
\[
 \ell(P)^\top x=r(P)^\top B(P)x=r(P)^\top b(P),
\]
which is independent of \(x\).
\end{proof}
\par\smallskip\noindent \textit{Justification.} The inherited lemma isolates observable algebra from causal validity. \textit{Gap.} ZhangLiMiaoTchetgen2023 already proves the abstract null-annihilator principle; this lemma only records its six-dimensional coupled WBC form. \textit{Consumer.} The headline theorem uses it to turn one latent-valid bridge into a value shared by all observed bridge solutions.

\begin{lemma}[lem:latent-bridge-identifies]\label{lem:latent-bridge-identifies}
For every \(\mathcal S\in\mathfrak M_{\mathrm{WBC}}^{+}\) and every \(x\in\mathcal H_{\mathrm{lat}}(\mathcal S)\cap\mathcal H_{\mathrm{obs}}(P_{\mathcal S}^O)\), the maintained WBC sequential, cross-world, consistency, and proxy assumptions, together with the two proposal-specific target-relevant latent transport-overlap restrictions, imply \(\psi(\mathcal S)=\ell(P_{\mathcal S}^O)^{\top}x\). The observable bridge system requires only \(P_{\mathcal S}^O(A=0,Z=z,M=m)>0\) and \(P_{\mathcal S}^O(A=1,Z=z)>0\) at its displayed indices; no completeness, injectivity, or positivity of individual \(W\) or \(Y\) values is used.
\end{lemma}
\begin{proof}
Put \(P=P_{\mathcal S}\), and let \((h_0,h_1)\) be the pair represented by
\(x\).  For \(h\in\mathcal H\) and \(m\in\mathcal X\), define
\[
 p_h=P(H_0=h),\qquad s_h(m)=P(M_{1,h}=m),\qquad
 g_h(m)=\mathbb E_P[Y_{0,h,m,W_h}\mid H_0=h]
\]
whenever \(p_hs_h(m)>0\).  The conditional expectations used below are well
defined: \(\mathrm{ass{:}latent\text{-}a0\text{-}transport\text{-}support}\)
gives \(P(H=h,M=m,A=0)>0\) for such \((h,m)\), and
\(\mathrm{ass{:}latent\text{-}a1\text{-}transport\text{-}support}\) gives
\(P(H=h,A=1)>0\) whenever \(p_h>0\).  The displayed observed bridge
conditionals are well defined by
\(\mathrm{ass{:}a0\text{-}bridge\text{-}denominators}\) and
\(\mathrm{ass{:}a1\text{-}bridge\text{-}denominators}\).
In particular, summing any positive displayed denominator cell over its remaining
indices gives \(P(A=0)>0\) and \(P(A=1)>0\), so every arm conditional below is
well defined.

We first connect \(g_h(m)\) to the first latent bridge.  On
\(\{H=h,M=m,A=0\}\), the consistency assumptions give
\(H=H_0\), \(W=W_h\), and \(Y=Y_{0,h,m,W_h}\).  Hence
\[
 \mathbb E_P[Y\mid H=h,M=m,A=0]
 =\mathbb E_P[Y_{0,h,m,W_h}\mid H=h,M=m,A=0].
\]
By \(\mathrm{ass{:}sequential\text{-}outcome}\), the right side is unchanged
when \(M=m\) is removed from the conditioning set.  By
\(\mathrm{ass{:}treatment\text{-}consistency}\) and
\(\mathrm{ass{:}sequential\text{-}treatment}\), it is then unchanged when
\(A=0\) is removed after rewriting \(H=h\) as \(H_0=h\).  Therefore
\[
 \mathbb E_P[Y\mid H=h,M=m,A=0]=g_h(m).
\]
Since \(x\in\mathcal H_{\mathrm{lat}}(\mathcal S)\), the first equation in
\(\mathrm{def{:}latent\text{-}bridge\text{-}set}\) and
\(\mathrm{ass{:}w\text{-}proxy\text{-}exclusion}\) yield
\[
 g_h(m)=\sum_{w\in\mathcal X}h_0(w,m)P(W=w\mid H=h).
 \tag{1}
\]
The same exclusion justifies replacing the displayed conditional law by its
versions conditional on \((M,A)\) wherever those versions occur.

Next, \(\mathrm{ass{:}sequential\text{-}mediator}\),
\(\mathrm{ass{:}treatment\text{-}consistency}\), and
\(\mathrm{ass{:}m\text{-}consistency}\) give, for every target-relevant \(h\),
\[
 s_h(m)=P(M_{1,h}=m\mid H=h,A=1)=P(M=m\mid H=h,A=1).
 \tag{2}
\]
Combining (1), (2), and
\(\mathrm{ass{:}w\text{-}proxy\text{-}exclusion}\), and then using the second
latent bridge equation, gives
\[
 \begin{aligned}
 \sum_m s_h(m)g_h(m)
 &=\mathbb E_P[h_0(W,M)\mid H=h,A=1]\\
 &=\mathbb E_P[h_1(W)\mid H=h,A=1]\\
 &=\mathbb E_P[h_1(W)\mid H=h,A=0].
 \end{aligned}
 \tag{3}
\]

It remains to preserve the same unit-level hidden witness in the nested target.
For each fixed \((h,m)\), \(\mathrm{ass{:}cross\text{-}world}\) states that
\(\{Y_{0,H_0,m,W_{H_0}},H_0\}\) is independent of \(M_{1,h}\).  Conditioning
on \(H_0=h\) and summing over the finite supports therefore gives
\[
 \psi(\mathcal S)
 =\sum_{h\in\mathcal H}p_h\sum_{m\in\mathcal X}s_h(m)g_h(m).
 \tag{4}
\]
Moreover, \(H_0\perp A\) is a marginal consequence of
\(\mathrm{ass{:}sequential\text{-}treatment}\), and
\(H=H_0\) on \(A=0\) by treatment consistency.  Thus
\(p_h=P(H=h\mid A=0)\).  Substitution of (3) into (4) yields
\[
 \psi(\mathcal S)=\sum_hP(H=h\mid A=0)
 \mathbb E_P[h_1(W)\mid H=h,A=0]
 =\mathbb E_P[h_1(W)\mid A=0].
\]
Finally, the definition of \(\ell(P_{\mathcal S}^O)\) makes the last expression
equal to \(\ell(P_{\mathcal S}^O)^\top x\).  No inversion, completeness, or
positivity of a value of \(W\) or \(Y\) entered the derivation.
\end{proof}
\par\smallskip\noindent \textit{Justification.} This is the causal step that keeps the same unit-level hidden witness through both nests. \textit{Gap.} WuBaiCui2026 Proposition 3.1 derives latent bridge equations by completeness; this lemma takes primitive latent-valid realizability and proves the nested formula without completeness. \textit{Consumer.} Binary-proxy mediation analyses can justify the nested mean even when the observed bridge system has many solutions.

\begin{proposition}[prop:positive-rank-two-success]\label{prop:positive-rank-two-success}
The exact SCM \(\mathcal S^+\) belongs to both \(\mathfrak M_{\mathrm{cal}}\) and \(\mathfrak M_{\mathrm{WBC}}^{+}\): it is the calibrated construction with \(q=(2/3,1/6,2/3)\), \(\alpha=1/2\), \(\tau=\beta=d=0\), and \(c=1/2\). It induces \(P^+(o)=1/32\) for all observed cells, has both three-state proxy-completeness maps of rank exactly \(2\) with explicit null vectors, and has \(\operatorname{rank}(B(P^+))=3<6\). Nevertheless \(\mathsf{RI}(P^+)\) holds, two distinct observed bridge solutions have the same loading, and \(\psi(\mathcal S^+)=1/2\).
\end{proposition}
\begin{proof}
Let the structural equations be those in \(\mathrm{def{:}success\text{-}artifact}\).  First consider the calibrated parameter tuple
\[
 q=(2/3,1/6,2/3),\qquad \alpha=1/2,\qquad c=1/2,
 \qquad \tau=\beta=d=0.
\]
For every \(h\in\mathcal H\), \(c+dq_h=1/2\in[0,1]\), and for every
\((m,w)\in\mathcal X^2\), \(\alpha+\tau m+\beta w=1/2\in[0,1]\).
Thus the tuple belongs to \(\Theta_{\mathrm{cal}}\).  The displayed structural
equations give
\[
 W_h=\mathbf 1\{U_W\le q_h\},\qquad
 M_{1,h}=\mathbf 1\{U_M\le c+dq_h\},\qquad
 Y_{0,h,m,w}=\mathbf 1\{U_Y\le\alpha+\tau m+\beta w\},
\]
so the artifact is generated by \(\mathrm{def{:}calibrated\text{-}bridge\text{-}family}\).

We verify the member properties rather than treating the construction as their
proof.  Consistency holds because each observed variable is obtained by evaluating
its displayed response function at its observed parents.  The mutually independent
shocks separate \(U_A\), \(U_Z\), \(U_W\), \(U_M\), and \(U_Y\).  Hence treatment is
independent of the relevant potential outcomes, every \(M_{a,h}\) is independent of
\((H,A)\), the cross-world mediator independence holds, and the fixed-\(m\) outcome
is conditionally independent of the observed mediator given \((H,A)\).  The separate
\(U_Z\) and \(U_W\) shocks give the two proxy exclusions.  Moreover,
\[
 P(H_0=h)=\frac13,\quad P(M_{1,h}=m)=\frac12,\quad
 P(A=0,H=h,M=m)=\frac1{12},\quad P(A=1,H=h)=\frac16.
\]
All antecedent-positive cases in the two latent transport-overlap conditions
therefore have positive required denominators.  The observed denominator conditions
will also follow from the strictly positive observed law computed below.  These
checks prove \(\mathcal S^+\in\mathfrak M_{\mathrm{cal}}\).

For latent bridge existence, take
\[
 h_0(w,m)=\frac12,\quad h_1(w)=\frac12.
\]
Since every outcome risk is \(1/2\), for all \((h,m)\),
\[
 \mathbb E[Y\mid H=h,M=m,A=0]=\frac12
 =\sum_w h_0(w,m)P(W=w\mid H=h,M=m,A=0).
\]
Similarly, for every \(h\), both sides of the second latent bridge equation equal
\(1/2\).  Thus \(x^{\mathrm{lat}}=(1/2,1/2,1/2,1/2,1/2,1/2)^\top\) is latent-valid.

We next compute the observed law.  Conditional on \(H=h\), the binary proxy risks
are \(q_h^Z=(1/4,1/2,3/4)_h\) and
\(q_h^W=(2/3,1/6,2/3)_h\).  Their means are both \(1/2\), and
\[
 \frac13\sum_{h=0}^2q_h^Zq_h^W
 =\frac13\left(\frac16+\frac1{12}+\frac12\right)=\frac14.
\]
Because the two proxy shocks are independent conditional on \(H\), these identities
give \(P(Z=z,W=w)=1/4\) for every \((z,w)\).  The treatment, mediator, and outcome
are independent fair bits under this artifact, so
\[
 P^+(a,z,w,m,y)=\frac12\cdot\frac14\cdot\frac12\cdot\frac12
 =\frac1{32}
\]
for every observed cell.  In particular, every denominator required by the bridge
system is positive.  Direct substitution into its six conditional moments gives
the displayed \(B(P^+)\), \(b(P^+)\), and \(\ell(P^+)\).  Multiplication gives
\(B(P^+)x^{\mathrm{lat}}=b(P^+)\), so the latent-valid vector also lies in the
observed bridge fiber.  Consequently latent-valid bridge existence holds, and the
remaining verified properties prove
\(\mathcal S^+\in\mathfrak M_{\mathrm{WBC}}^+\).

The two conditional-proxy matrices can be checked exactly.  With rows ordered by
proxy value \(0,1\),
\[
 P(H=\cdot\mid Z=\cdot)=
 \begin{pmatrix}1/2&1/3&1/6\\[2pt]1/6&1/3&1/2\end{pmatrix},
 \qquad
 P(H=\cdot\mid W=\cdot)=
 \begin{pmatrix}2/9&5/9&2/9\\[2pt]4/9&1/9&4/9\end{pmatrix}.
\]
The first matrix annihilates \((1,-2,1)^\top\), and the second annihilates
\((1,0,-1)^\top\).  In each matrix the two rows are unequal, so each has rank
exactly \(2\), not \(3\).  Conditioning additionally on the independent treatment
or mediator variables leaves these ranks unchanged.

For the coupled matrix, its first, third, and fifth displayed rows are linearly
independent and every other row duplicates one of them.  Hence
\(\operatorname{rank}(B(P^+))=3\).  Equivalently, its null space has the displayed
basis \(e_1-e_2,e_3-e_4,e_5-e_6\).  The vector
\[
 r=(1/2,0,1/2,0,-1,0)^\top
\]
satisfies \(B(P^+)^\top r=\ell(P^+)\), while
\(x^{(1)}=(1,0,1,0,1,0)^\top\) and
\(x^{(2)}=(0,1,0,1,0,1)^\top\) both satisfy \(B(P^+)x=b(P^+)\).  Therefore
\(\mathsf{RI}(P^+)\) holds, and
\[
 \ell(P^+)^\top x^{(1)}=\ell(P^+)^\top x^{(2)}=\frac12.
\]
Finally, the nested outcome is a Bernoulli threshold with success probability
\(1/2\) for every value of all its arguments, so
\(\psi(\mathcal S^+)=1/2\).  This proves every assertion of the proposition.
\end{proof}
\par\smallskip\noindent \textit{Justification.} The rational SCM makes the below-completeness identifying branch nonvacuous. \textit{Gap.} Neither MiaoGengTchetgen2018, WuBaiCui2026, ShpitserWoodDoughtyTchetgen2023, nor BaiWuSunCui2026 supplies a positive three-state WBC SCM inside a primitive calibrated family with two rank-two binary channels, a rank-three coupled operator, and scalar target invariance; Bai--Wu--Sun--Cui instead study an observed recanting witness distinct from the latent confounder under three stage-specific completeness conditions. \textit{Consumer.} It is an auditable benchmark for the NLSY97 analysis after binary coarsening of each behavioral proxy block.

\begin{theorem}[thm:calibrated-target-span-identification]\label{thm:calibrated-target-span-identification}
For every \(\mathcal S^{\mathrm{cal}}\in\mathfrak M_{\mathrm{cal}}\), \(h_0(w,m)=\alpha+\tau m+\beta w\) and \(h_1(w)=\alpha+\tau c+(\beta+\tau d)w\) obey \(\mathbb E[Y\mid H=h,M=m,A=0]=\mathbb E[h_0(W,m)\mid H=h,M=m,A=0]\) at every target-relevant \((h,m)\), \(\mathbb E[h_0(W,M)\mid H=h,A=1]=\mathbb E[h_1(W)\mid H=h,A=1]\) at every target-relevant \(h\), and \(\psi(\mathcal S^{\mathrm{cal}})=\mathbb E[h_1(W)\mid A=0]\). The same pair satisfies the six observed equations, so \(\mathcal H_{\mathrm{lat}}(\mathcal S^{\mathrm{cal}})\cap\mathcal H_{\mathrm{obs}}(P_{\mathcal S^{\mathrm{cal}}}^O)\neq\varnothing\), observed bridge solvability follows, and \(\mathfrak M_{\mathrm{cal}}\subseteq\mathfrak M_{\mathrm{WBC}}^{+}\). If \(\ell(P_{\mathcal S^{\mathrm{cal}}}^O)\in\operatorname{row}(B(P_{\mathcal S^{\mathrm{cal}}}^O))\), every observed bridge solution has common value \(\psi(\mathcal S^{\mathrm{cal}})=r(P_{\mathcal S^{\mathrm{cal}}}^O)^{\top}b(P_{\mathcal S^{\mathrm{cal}}}^O)\). The exact \(\mathcal S^+\) is a member, so this sufficient subclass identifies a positive WBC law at which both three-state proxy-completeness conditions fail. This theorem and Shpitser--Wood-Doughty--Tchetgen Tchetgen Theorem 4 do not subsume one another: their Assumptions 2.5 and 5.13 retain injectivity at the relevant proxy step and their conclusion is a recursive interventional margin, whereas this result concerns one scalar same-unit hidden-recanting mean under WBC-specific causal and latent transport restrictions and gives no general proximal-fixing algorithm. This theorem and Bai--Wu--Sun--Cui Theorem 2.1 also do not subsume one another. Their Assumption 6(i)--(iii) retains three stage-specific completeness conditions for a latent baseline confounder, their recanting witness is observed and distinct from that confounder, and their target and three-bridge construction belong to a different path-specific graph. Conversely, this result treats only the fixed hidden-witness WBC graph and its scalar same-unit target, relies on the calibrated subclass or the explicit latent-valid intersection and two non-testable latent transport restrictions, and supplies neither their observed-witness construction nor their inference results.
\end{theorem}
\begin{proof}
Fix \(\mathcal S^{\mathrm{cal}}\in\mathfrak M_{\mathrm{cal}}\), and write
\(P=P_{\mathcal S^{\mathrm{cal}}}^O\).  Membership in the calibrated class supplies
the common-shock equations in \(\mathrm{def{:}calibrated\text{-}bridge\text{-}family}\)
and all member properties listed in \(\mathrm{def{:}calibrated\text{-}wbc\text{-}subclass}\).
For each \(h\), independence of \(U_W\) from the remaining exogenous variables gives
\[
 P(W_h=1)=q_h,
\]
and independence of \(U_M\) gives
\[
 P(M_{1,h}=1)=c+dq_h.
\]
The inequalities defining \(\Theta_{\mathrm{cal}}\) make these valid probabilities
and make every threshold \(\alpha+\tau m+\beta w\) lie in \([0,1]\).

Define
\[
 h_0(w,m)=\alpha+\tau m+\beta w,
 \qquad
 h_1(w)=\alpha+\tau c+(\beta+\tau d)w.
 \tag{1}
\]
For every \((h,m)\) with \(P(H=h,M=m,A=0)>0\), consistency, the arm-zero
outcome equation, and the independence of \(U_Y\) and \(U_W\) give
\[
 \begin{aligned}
 \mathbb E[Y\mid H=h,M=m,A=0]
 &=\alpha+\tau m+\beta q_h\\
 &=\sum_{w\in\mathcal X}(\alpha+\tau m+\beta w)
          P(W=w\mid H=h,M=m,A=0)\\
 &=\mathbb E[h_0(W,m)\mid H=h,M=m,A=0].
 \end{aligned}
 \tag{2}
\]
No division by a proxy-cell probability occurs in (2).  In particular, for every
target-relevant \((h,m)\), the arm-zero latent transport support supplies the positive
conditioning event required in (2); and for the later observed-law averaging step,
(2) already holds at every positive arm-zero latent cell.

For every \(h\) with \(P(H=h,A=1)>0\), consistency and the independent common
shocks yield
\[
 \begin{aligned}
 \mathbb E[h_0(W,M)\mid H=h,A=1]
 &=\alpha+\tau(c+dq_h)+\beta q_h\\
 &=\alpha+\tau c+(\beta+\tau d)q_h\\
 &=\mathbb E[h_1(W)\mid H=h,A=1].
 \end{aligned}
 \tag{3}
\]
For each target-relevant \(h\), the arm-one latent transport support supplies the
positive conditioning event required in (3).  Equations (2)--(3) are therefore the
two latent bridge equations at every index required by
\(\mathrm{def{:}latent\text{-}bridge\text{-}set}\), so the vector induced by (1)
belongs to \(\mathcal H_{\mathrm{lat}}(\mathcal S^{\mathrm{cal}})\).  Their validity
on every positive observed latent cell will also justify the tower-property averages
below.

We next derive the causal equality rather than defining the causal target by the
bridge.  Let \(\pi_0(h)=P(H_0=h)\).  Conditional on \(H_0=h\), independence of
\(U_Y,U_M,U_W\) gives
\[
 \begin{aligned}
 \mathbb E[Y_{0,h,M_{1,h},W_h}\mid H_0=h]
 &=\alpha+\tau\mathbb E[M_{1,h}]+\beta\mathbb E[W_h]\\
 &=\alpha+\tau(c+dq_h)+\beta q_h.
 \end{aligned}
 \tag{4}
\]
Therefore, by the definition of \(\psi\),
\[
 \psi(\mathcal S^{\mathrm{cal}})
 =\sum_{h\in\mathcal H}\pi_0(h)
   \{\alpha+\tau c+(\beta+\tau d)q_h\}.
 \tag{5}
\]
Treatment consistency and sequential treatment independence imply that the law of
\(H\) conditional on \(A=0\) is the law of \(H_0\).  The outcome-proxy exclusion,
or directly the independent \(U_W\) construction, gives
\(P(W=1\mid H=h,A=0)=q_h\).  Hence
\[
 \mathbb E[h_1(W)\mid A=0]
 =\sum_h\pi_0(h)\{\alpha+\tau c+(\beta+\tau d)q_h\}
 =\psi(\mathcal S^{\mathrm{cal}}).
 \tag{6}
\]

It remains to prove that the same bridge is observable.  The two observable
denominator assumptions ensure that every conditional moment used below is defined.
By the treatment-inducing proxy exclusion and the tower property for finite sums,
\[
 \mathbb E[Y\mid Z=z,M=m,A=0]
 =\mathbb E\!\left[\mathbb E[Y\mid H,M=m,A=0]
                    \mid Z=z,M=m,A=0\right].
\]
Insert (2).  The outcome-inducing proxy exclusion implies that the conditional law
of \(W\) given \((H,Z,M,A)\) equals its conditional law given \(H\).  Applying the
tower property again gives
\[
 \mathbb E[Y\mid Z=z,M=m,A=0]
 =\mathbb E[h_0(W,m)\mid Z=z,M=m,A=0].
 \tag{7}
\]
For arm one, the two proxy exclusions jointly factor the conditional law of
\((W,M)\) given \((H,Z,A=1)\) as the law of \(W\) given \(H\) times the law of
\(M\) given \((H,A=1)\).  Thus (3) and two applications of the tower property give
\[
 \mathbb E[h_0(W,M)\mid Z=z,A=1]
 =\mathbb E[h_1(W)\mid Z=z,A=1].
 \tag{8}
\]
By \(\mathrm{def{:}coupled\text{-}bridge\text{-}matrix}\), equations (7)--(8) are
precisely the six equations \(B(P)x=b(P)\).  Thus
\[
 x\in\mathcal H_{\mathrm{lat}}(\mathcal S^{\mathrm{cal}})
       \cap\mathcal H_{\mathrm{obs}}(P).
 \tag{9}
\]
The intersection-implies-solvability lemma now gives observed solvability.  The
calibrated class already imposes every other member property of
\(\mathfrak M_{\mathrm{WBC}}^+\), and (9) verifies its remaining latent-valid bridge
existence property.  Therefore
\[
 \mathfrak M_{\mathrm{cal}}\subseteq\mathfrak M_{\mathrm{WBC}}^+.
 \tag{10}
\]
The latent-bridge identification lemma is applicable to (9) and is consistent with
the direct calculation (6):
\[
 \psi(\mathcal S^{\mathrm{cal}})=\ell(P)^\top x.
 \tag{11}
\]

Suppose now that \(\ell(P)\in\operatorname{row}(B(P))\).  By (9),
\(b(P)\in\operatorname{col}(B(P))\).  The null-annihilator lemma therefore implies
that \(\ell(P)^\top\widetilde x\) is the same for every
\(\widetilde x\in\mathcal H_{\mathrm{obs}}(P)\).  If
\(B(P)^\top r(P)=\ell(P)\), then for every such \(\widetilde x\),
\[
 \ell(P)^\top\widetilde x
 =r(P)^\top B(P)\widetilde x
 =r(P)^\top b(P).
 \tag{12}
\]
Taking \(\widetilde x=x\) in (12) and using (11) proves the asserted row-span
identifying formula.  This argument uses no completeness or injectivity and no
positivity for individual values of \(W\) or \(Y\).

The positive-rank-two-success proposition proves that the exact
\(\mathcal S^+\) is a member of the calibrated class, has target \(1/2\), and has
rank-two rather than rank-three proxy maps.  Hence the subclass is nonempty at a
strictly positive observed law where both stated three-state completeness conditions
fail.  Finally, the cited Proximal ID scope lemma records that Shpitser--Wood-Doughty--
Tchetgen Tchetgen retain ordinary and sequential injectivity and identify a recursive
interventional margin, while the cited proximal path-specific scope lemma records that
Bai--Wu--Sun--Cui use an observed recanting witness distinct from a latent baseline
confounder, three completeness conditions, and a different three-bridge target.
Equations (1)--(12) concern instead the fixed hidden-witness WBC graph and its scalar
same-unit nested mean, and they provide neither cited paper's general construction or
inference conclusions.  Conversely, the cited conclusions do not supply (1)--(12) or
the rank-deficient member just established.  These differences prove the two stated
non-subsumption comparisons without enlarging the scope of the present theorem.
\end{proof}
\par\smallskip\noindent \textit{Justification.} This is the principal substantive result: an interpretable primitive family derives bridge realizability and target identification without completeness. \textit{Gap.} WuBaiCui2026 obtains latent bridge validity through completeness. ShpitserWoodDoughtyTchetgen2023 Assumptions 2.5 and 5.13 likewise impose ordinary and sequential latent injectivity, while its Theorem 4 identifies the next recursive interventional margin under sequential proxy independence, fixing ignorability, bridge existence, completeness, and positivity. BaiWuSunCui2026 Section 2 and Figure 1(c), Assumption 5, Assumption 6(i)--(iii), and Theorem 2.1, equations (4)--(7), instead treat an observed treatment-induced recanting witness distinct from the latent baseline confounder and use three stage-specific completeness conditions to identify their different path-specific mean. This theorem identifies the fixed WBC same-unit scalar at a positive law where both three-state binary-proxy completeness maps fail; it neither covers nor is covered by the Bai--Wu--Sun--Cui graph, target, bridge sequence, or inference results. \textit{Consumer.} NLSY97 analysts can justify the high-school-track-to-wage nested mean through calibrated latent severity, affine achievement risk, and additive wage effects, subject to the explicit non-testable transport restrictions.

\begin{theorem}[thm:wbc-target-span-identification]\label{thm:wbc-target-span-identification}
For every \(\mathcal S\in\mathfrak M_{\mathrm{WBC}}^{+}\), put \(P=P_{\mathcal S}^O\). The class premise \(\mathcal H_{\mathrm{lat}}(\mathcal S)\cap\mathcal H_{\mathrm{obs}}(P)\neq\varnothing\) implies \(b(P)\in\operatorname{col}(B(P))\). If \(\ell(P)\in\operatorname{row}(B(P))\), every \(x\in\mathcal H_{\mathrm{obs}}(P)\) satisfies \(\psi(\mathcal S)=\ell(P)^{\top}x=r(P)^{\top}b(P)\). Hence \(\psi\) is point identified over all members of \(\mathfrak M_{\mathrm{WBC}}^{+}\) inducing \(P\), and on observed laws induced by this class \(\mathsf{RI}(P)\) is equivalent to the single loading condition \(\ell(P)\in\operatorname{row}(B(P))\). This is a conditional envelope rather than an observable or primitive frontier: it assumes nonempty latent-valid observed bridge intersection and the two proposal-specific target-relevant latent transport-overlap restrictions. It permits zero probabilities for values of \(W\) and \(Y\) and imposes only the displayed bridge denominators on the observed law.
\end{theorem}
\begin{proof}
Fix \(\mathcal S\in\mathfrak M_{\mathrm{WBC}}^+\) and put
\(P=P_{\mathcal S}^O\).  The class definition includes
\(\mathrm{ass{:}latent\text{-}valid\text{-}bridge\text{-}existence}\), so
\(\mathrm{lem{:}intersection\text{-}implies\text{-}solvability}\) gives
\(b(P)\in\operatorname{col}(B(P))\).  Suppose now that
\(\ell(P)\in\operatorname{row}(B(P))\).  Choose one
\(x_\star\in\mathcal H_{\mathrm{lat}}(\mathcal S)\cap
\mathcal H_{\mathrm{obs}}(P)\).  By
\(\mathrm{lem{:}latent\text{-}bridge\text{-}identifies}\),
\[
 \psi(\mathcal S)=\ell(P)^\top x_\star.
\]
The row-space condition and \(\mathrm{lem{:}null\text{-}annihilator}\) make the
loading constant over the entire observed affine fiber, and for every representer
\(B(P)^\top r(P)=\ell(P)\) give
\[
 \ell(P)^\top x=\ell(P)^\top x_\star
 =\psi(\mathcal S)=r(P)^\top b(P),
 \qquad x\in\mathcal H_{\mathrm{obs}}(P).
\]
The rightmost expression is a function only of the observed law.  Therefore any
two class members inducing the same \(P\) have the same target value, which is
point identification rather than a definition of the causal target.

On a law induced by this class, observed solvability has already been proved.
Thus the column-space conjunct in
\(\mathrm{def{:}row\text{-}space\text{-}criterion}\) is automatic, and
\(\mathsf{RI}(P)\) is equivalent to the remaining row-space condition.  The only
positivity conditions used to form the observed system are the two displayed
denominator assumptions; neither the proof nor the definitions divide by a
probability of a particular value of \(W\) or \(Y\).
\end{proof}
\par\smallskip\noindent \textit{Justification.} The theorem records the broadest honest conditional envelope for the fixed hidden-recanting mean without presenting latent bridge validity as observable. \textit{Gap.} WuBaiCui2026 uses completeness for the causal lift; ZhangLiMiaoTchetgen2023 gives abstract one-stage invariance but not this shared-witness causal equality. The primitive content establishing that the envelope is usable is supplied by the calibrated-subclass theorem. \textit{Consumer.} It identifies the NLSY97 wage return operating through cognitive achievement while the same latent non-cognitive trait remains at its general-track value.

\begin{proposition}[prop:scm-counterpair]\label{prop:scm-counterpair}
The exact SCMs \(\mathcal S^{\oplus}\) and \(\mathcal S^{\ominus}\) belong to \(\mathfrak M_{\mathrm{WBC}}^{+}\), induce the same strictly positive \(P^{\dagger}\), and have the same displayed \(B(P^{\dagger})\), \(b(P^{\dagger})\), and \(\ell(P^{\dagger})\). The row-space criterion fails, and \(\psi(\mathcal S^{\oplus})=19/48\neq7/16=\psi(\mathcal S^{\ominus})\). This is causal nonidentification at one positive law, not a universal necessity theorem.
\end{proposition}
\begin{proof}
Put
\[
 \pi _0=(1/2,1/3,1/6),\qquad \pi _1=(1/3,1/3,1/3),
 \qquad q=(1/4,1/2,3/4),
\]
and write
\[
 s^+=(5/8,1/4,5/8),\qquad s^-=(3/8,3/4,3/8),
 \qquad t=(0,1/2,1).
\]
All calculations below are exact rational calculations.

First we verify the structural and independence premises in
\(\mathrm{def{:}wbc\text{-}class}\).  In either model, each potential response is
a measurable function of the independent shock specified in
\(\mathrm{def{:}failure\text{-}artifact}\): \(A\) uses \(U_A\), the entire
response \((H_0,H_1)\) uses \(U_H\), all \(Z_{a,h}\) use \(U_Z\), all
\(W_h\) use \(U_W\), all \(M_{a,h}\) use \(U_M\), and all
\(Y_{a,h,m,w}\) use \(U_Y\).  The pair \((U_H,U_Y)\) may be regarded as
the exogenous block for the district \(\{H,Y\}\); its two coordinates happen
to be independent, which is allowed.  The displayed structural substitutions
therefore have precisely the parent sets allowed by
\(\mathcal G_{\mathrm{WBC}}\), and the five consistency equalities hold by
construction.

Moreover, \(A\), being a function of \(U_A\), is independent of every
variable in
\(\{Y_{a,H_a,m,W_{H_a}},H_a\}\).  The variable \(M_{a',h}\) is a
function of \(U_M\), whereas
\(\{Y_{a,H_a,m,W_{H_a}},H_a\}\) is a function of
\((U_H,U_W,U_Y)\); this proves the cross-world independence.  The same
block independence gives \(M_{a,h}\mathbin{\perp\!\!\!\perp}(H,A)\).
Conditional on \((H,A)\), the realized mediator is a function of \(U_M\),
whereas \(Y_{A,H,m,W_H}\) is a function of \((U_W,U_Y)\); hence the
sequential-outcome independence holds.  Conditional on \((H,A)\), \(Z\)
uses only the independent shock \(U_Z\), proving
\(Z\mathbin{\perp\!\!\!\perp}(Y,M)\mid(H,A)\).  Conditional on \(H\),
\(W\) uses only \(U_W\), independently of the shocks determining
\((A,Z,M)\), proving
\(W\mathbin{\perp\!\!\!\perp}(A,Z,M)\mid H\).  Thus every listed
sequential, cross-world, consistency, and proxy-exclusion premise is
satisfied.

The observable denominator probabilities are
\[
 P(A=0,Z=z,M=m)=\frac12\frac12\frac12=\frac18>0,
 \qquad
 P(A=1,Z=z)=\frac12\frac12=\frac14>0.
\]
Every coordinate of \(\pi _0,\pi _1,s^+,s^-\) lies strictly between zero
and one.  Consequently, for every \(h,m\),
\[
 P(H_0=h)>0,\quad P(M_{1,h}=m)>0,\quad
 P(H=h,M=m,A=0)=\frac14\pi _0(h)>0,\quad
 P(H=h,A=1)=\frac16>0.
\]
These equations verify both target-relevant latent transport-overlap
restrictions.

It remains to verify latent-valid observed bridge existence.  In both models
take
\[
 h_0(0,0)=h_0(1,0)=\frac12,\qquad
 h_0(0,1)=-\frac12,\qquad h_0(1,1)=\frac32.
\]
Since \(W_h\) is Bernoulli with mean \(q_h\),
\[
 E\{h_0(W_h,0)\}=\frac12,\qquad
 E\{h_0(W_h,1)\}=-\frac12+2q_h=t_h.
\]
These are exactly the treatment-zero conditional outcome means for mediator
values zero and one, respectively, so the first latent bridge equations hold.
For the second bridge, conditional independence of \(W_h\) and
\(M_{1,h}\) yields
\[
 L_h^{\pm}:=E\{h_0(W_h,M_{1,h})\}
   =(1-s_h^{\pm})\frac12+s_h^{\pm}t_h.
\]
Direct substitution gives
\[
 (L_0^+,L_1^+,L_2^+)=(3/16,1/2,13/16),\qquad
 (L_0^-,L_1^-,L_2^-)=(5/16,1/2,11/16).
\]
For the two proposed second-stage bridges,
\[
\begin{aligned}
 E\{h_1^+(W_h)\}
   &=-\frac18(1-q_h)+\frac98q_h=L_h^+,\\
 E\{h_1^-(W_h)\}
   &= \frac18(1-q_h)+\frac78q_h=L_h^-.
\end{aligned}
\]
Thus the displayed vectors \(x^+\) and \(x^-\) in
\(\mathrm{def{:}failure\text{-}artifact}\) belong to their respective
latent bridge sets.  They also solve the observable equations.  Indeed, for
either vector the first two distinct row calculations are
\[
 \frac7{12}\frac12+\frac5{12}\frac12=\frac12,
 \qquad
 \frac7{12}\left(-\frac12\right)
   +\frac5{12}\frac32=\frac13,
\]
and the treatment-one row is
\[
 \frac14\left(\frac12+\frac12-\frac12+\frac32\right)
 -\frac12\{h_1(0)+h_1(1)\}=\frac12-\frac12=0,
\]
because the two coordinates of either \(h_1\) sum to one.  Hence
\(B(P^{\dagger})x^{\pm}=b(P^{\dagger})\), proving the required nonempty
latent-valid observed bridge intersection.  We have now verified every member
property of \(\mathfrak M_{\mathrm{WBC}}^{+}\).

We next verify observational equivalence.  The two models are identical in
the treatment-zero arm.  In the treatment-one arm, \(Z\) and \(Y\) are
independent Bernoulli variables with mean \(1/2\).  For either sign,
\[
 \frac13\sum_h s_h^{\pm}=\frac12,\qquad
 \frac13\sum_h q_h=\frac12,\qquad
 \frac13\sum_h q_hs_h^{\pm}=\frac14.
\]
Because \(W\) and \(M\) are binary, these three moments imply
\(P(W=w,M=m\mid A=1)=1/4\) for every \((w,m)\).  Thus the entire observed
law agrees and is the displayed \(P^{\dagger}\).  Every treatment-one cell
is positive.  In the treatment-zero arm, \(Z\) and \(M\) have probability
\(1/2\) at each value and every \(q_h\) is interior; for \(M=0\), both
outcome values have probability \(1/2\), while for \(M=1\) the positive
\(h=1\) mixture component assigns probability \(1/2\) to either outcome.
Therefore every cell of \(P^{\dagger}\) is strictly positive.

Under \(A=0\), \(M\) is independent of \(H\), and hence
\[
 P(W=1\mid A=0,Z=z,M=m)=\sum_h\pi _0(h)q_h=\frac5{12}.
\]
The treatment-zero outcome means are \(1/2\) for \(m=0\) and
\(\sum_h\pi _0(h)t_h=1/3\) for \(m=1\).  Under \(A=1\), the preceding
calculation gives \(P(W=w,M=m\mid Z=z,A=1)=1/4\) and
\(P(W=w\mid Z=z,A=1)=1/2\).  Substitution into
\(\mathrm{def{:}coupled\text{-}bridge\text{-}matrix}\) gives exactly the
displayed \(B(P^{\dagger})\), \(b(P^{\dagger})\), and
\(\ell(P^{\dagger})\).  The first, third, and fifth rows of the displayed
matrix are linearly independent and every other row duplicates one of these,
so its rank is three.  The vector
\[
 v=(0,0,0,0,1,-1)^{\top}
\]
satisfies \(B(P^{\dagger})v=0\), but
\[
 \ell(P^{\dagger})^{\top}v=\frac7{12}-\frac5{12}=\frac16\ne0.
\]
Every vector in the row space annihilates every vector in the null space;
therefore \(\ell(P^{\dagger})\notin
\operatorname{row}(B(P^{\dagger}))\), and the row-space criterion fails.

Finally, conditioning the nested target first on \(H_0=h\) gives
\[
 \psi(\mathcal S^{\pm})
 =\sum_{h=0}^2\pi _0(h)
 \left\{(1-s_h^{\pm})\frac12+s_h^{\pm}t_h\right\}.
\]
Consequently,
\[
\begin{aligned}
 \psi(\mathcal S^{\oplus})
 &=\frac12\frac3{16}+\frac13\frac12+\frac16\frac{13}{16}
   =\frac{19}{48},\\
 \psi(\mathcal S^{\ominus})
 &=\frac12\frac5{16}+\frac13\frac12+\frac16\frac{11}{16}
   =\frac7{16}.
\end{aligned}
\]
Their difference is \(1/24\), completing every assertion.
\end{proof}
\par\smallskip\noindent \textit{Justification.} The counterpair upgrades a bridge-vector ambiguity into genuine observationally equivalent causal models. \textit{Gap.} GhassamiZhangShpitserTchetgen2026 treats other incomplete-proxy targets; this proposition supplies an exact WBC counterpair rather than a comparison to any in-repository artifact. \textit{Consumer.} It is an auditable benchmark demonstrating genuine causal ambiguity at one fixed observed law.

\begin{theorem}[thm:minimal-binary-augmentation]\label{thm:minimal-binary-augmentation}
Fix the displayed rank-two \(Z\)-channel and allow added treatment-inducing proxy coordinates obeying WBC exclusion and conditional independence given \((H,A)\). Zero added binary coordinates cannot give rank three for three-state \(H\), because a \(2\times3\) channel has rank at most \(2\). The displayed one-bit product channel satisfies \(\operatorname{rank}(R_{Z^{\star}\mid H})=3\). For every strictly positive law of \(H\) given \((M,A=0)\) and given \(A=1\), Bayes row and column reweighting preserves rank, so the two WBC completeness clauses hold. Hence one added binary treatment-inducing proxy coordinate is the sharp minimum for uniform completeness restoration for this fixed input channel. If, in addition, there exist functions \(h_0:\mathcal X^2\to\mathbb R\) and \(h_1:\mathcal X\to\mathbb R\) such that, under the augmented SCM law and almost surely on the indicated conditioning events,
\[
 \mathbb E[Y\mid Z^{\star},M,A=0]
 =\sum_{w\in\mathcal X}h_0(w,M)\Pr(W=w\mid Z^{\star},M,A=0)
\]
and
\[
 \mathbb E[h_0(W,M)\mid Z^{\star},A=1]
 =\sum_{w\in\mathcal X}h_1(w)\Pr(W=w\mid Z^{\star},A=1),
\]
and every remaining causal, consistency, exclusion, and support or positivity premise of Wu--Bai--Cui Theorem 3.1 holds, then that theorem's identification conclusion applies. Full augmented-channel rank alone establishes neither those remaining premises nor bridge existence.
\end{theorem}
\begin{proof}
Write the original binary channel, with rows ordered by \(Z=0,1\) and columns by \(H=0,1,2\), as
\[
 R_{Z\mid H}=
 \begin{pmatrix}
  3/4&1/2&1/4\\
  1/4&1/2&3/4
 \end{pmatrix}.
\]
The determinant of its first two columns is
\[
 \det\!\begin{pmatrix}3/4&1/2\\1/4&1/2\end{pmatrix}
 =\frac{3}{8}-\frac{1}{8}=\frac14\ne0,
\]
so this channel has rank two. More generally, without an added coordinate a proxy with binary alphabet has a \(2\times3\) channel from the three states of \(H\), and linear algebra gives rank at most two.

We first establish the exact finite completeness criterion used at both stages. Let \(S\) be any finite-valued proxy, let \(C\) be a conditioning event with \(P(C)>0\), and define
\[
 R_{sh}=P(S=s\mid H=h,C),\qquad
 \pi_h=P(H=h\mid C).
\]
Assume \(\pi_h>0\) for every \(h\in\mathcal H\) and \((R\pi)_s>0\) for every supported proxy state \(s\). Bayes' rule gives
\[
 K_C(s,h):=P(H=h\mid S=s,C)
 =\frac{R_{sh}\pi_h}{(R\pi)_s},
 \qquad
 K_C=D_{R\pi}^{-1}R D_\pi,
\]
where \(D_\pi\) and \(D_{R\pi}\) are diagonal. The stated positivity makes both diagonal matrices invertible, whence
\[
 \operatorname{rank}(K_C)=\operatorname{rank}(R).
\]
For a real function \(g\) on the finite set \(\mathcal H\), the equations \(\mathbb E[g(H)\mid S,C]=0\) almost surely are precisely \(K_Cg=0\). Therefore full column rank of \(R\) implies \(g=0\). Conversely, if \(R\) has rank below three, then \(K_C\) has a nonzero null vector and completeness fails. Since \(\mathcal H\) is finite, every such \(g\) is square-integrable.

By the conditional independence imposed in \(\mathrm{def{:}binary\text{-}proxy\text{-}augmentation}\), for every \(a\in\mathcal A\) and \(h\in\mathcal H\),
\[
 P(Z=z,V=v\mid H=h,A=a)
 =P(Z=z\mid H=h,A=a)P(V=v\mid H=h,A=a).
\]
Multiplying the displayed Bernoulli probabilities gives
\[
 R_{Z^{\star}\mid H}=
 \begin{pmatrix}
  9/16&1/8&1/8\\
  3/16&3/8&1/8\\
  3/16&1/8&3/8\\
  1/16&3/8&3/8
 \end{pmatrix}.
\]
Its top-three-row minor equals
\[
 \frac{1}{16^3}
 \det\!\begin{pmatrix}9&2&2\\3&6&2\\3&2&6\end{pmatrix}
 =\frac{9(36-4)-2(18-6)+2(6-18)}{4096}
 =\frac{240}{4096}=\frac{15}{256}\ne0.
\]
Thus \(\operatorname{rank}(R_{Z^{\star}\mid H})=3\), the maximum possible column rank.

For the first WBC completeness clause, take \(C=\{A=1\}\). Assumption \(\mathrm{ass{:}treatment\text{-}positivity}\) gives \(P(C)>0\), and the theorem's strict-positivity hypothesis gives \(\pi_h=P(H=h\mid A=1)>0\) for every \(h\). Every entry of \(R_{Z^{\star}\mid H}\) is positive, so every entry of \(R_{Z^{\star}\mid H}\pi\) is positive. The preceding Bayes calculation yields
\[
 \mathbb E[g(H)\mid Z^{\star},A=1]=0
 \quad\Longrightarrow\quad
 g(H)=0\quad\text{almost surely}.
\]
For the second clause, fix \(m\in\mathcal X\) and take \(C=\{M=m,A=0\}\). Assumption \(\mathrm{ass{:}a0\text{-}bridge\text{-}denominators}\) implies \(P(C)>0\), since \(P(A=0,Z=z,M=m)>0\) for every \(z\). Assumption \(\mathrm{ass{:}z\text{-}proxy\text{-}exclusion}\), together with the conditional independence of \(V\) from \((Z,W,M,Y)\) given \((H,A)\) in \(\mathrm{def{:}binary\text{-}proxy\text{-}augmentation}\), gives
\[
 P(Z^{\star}=s\mid H=h,M=m,A=0)
 =P(Z^{\star}=s\mid H=h,A=0)
 =(R_{Z^{\star}\mid H})_{sh}.
\]
The theorem's strict-positivity hypothesis now supplies \(P(H=h\mid M=m,A=0)>0\) for every \(h\). Applying the same Bayes calculation gives
\[
 \mathbb E[g(H)\mid Z^{\star},M=m,A=0]=0
 \quad\Longrightarrow\quad
 g(H)=0\quad\text{almost surely}.
\]
This holds for every \(m\in\mathcal X\), so both WBC completeness clauses hold. With no added binary coordinate, full column rank three is impossible; with the single displayed binary coordinate, it holds for every one of the strictly positive latent laws specified in the theorem. Hence one added binary coordinate is necessary and sufficient, and is the sharp minimum for uniform completeness restoration for this fixed input channel.

The augmented proxy also inherits the required exclusions. In the finite model, \(\mathrm{ass{:}z\text{-}proxy\text{-}exclusion}\) and the conditional independence imposed on \(V\) imply
\[
 p(z,v,y,m\mid h,a)
 =p(v\mid h,a)p(z\mid h,a)p(y,m\mid h,a),
\]
which is \(Z^{\star}\mathbin{\perp\!\!\!\perp}(Y,M)\mid(H,A)\). Similarly, \(\mathrm{ass{:}w\text{-}proxy\text{-}exclusion}\) and the same conditional independence imply
\[
 p(w,a,z,v,m\mid h)
 =p(w\mid h)p(v\mid h,a)p(a,z,m\mid h)
 =p(w\mid h)p(a,z,v,m\mid h),
\]
which is \(W\mathbin{\perp\!\!\!\perp}(A,Z^{\star},M)\mid H\).

Now suppose that the two augmented bridge equations displayed in the theorem hold. They are the finite, no-baseline-covariate bridge equations in Wu--Bai--Cui Assumption 3.3 with treatment-inducing proxy \(Z^{\star}\). The two completeness clauses were proved above, and the augmented proxy exclusions follow from the displayed factorizations. Under the theorem's further hypothesis that every remaining causal, consistency, exclusion, support, and positivity premise of Wu--Bai--Cui Theorem 3.1 holds, all premises of \(\mathrm{lem{:}wbc\text{-}por\text{-}identification}\) are therefore satisfied. That cited result gives the explicit identification formula
\[
 \psi(\mathcal S)
 =\sum_{w\in\mathcal X}h_1(w)P_{\mathcal S}^O(W=w\mid A=0).
\]

It remains to prove that channel rank alone supplies neither the remaining premises nor bridge existence. For the first assertion, retain the displayed full-rank channel and set \(M=0\) almost surely. Then
\[
 P_{\mathcal S}^O(A=0,Z=z,M=1)=0
\]
for both \(z\in\mathcal X\), so \(\mathrm{ass{:}a0\text{-}bridge\text{-}denominators}\) fails even though the augmented channel still has rank three.

For bridge existence, consider the following explicit finite SCM. Let \(A\), \(M\), and \(W\) be mutually independent fair Bernoulli variables, let \(H\) be uniform on \(\mathcal H\), and take all four variables mutually independent. Generate \(Z\) and \(V\), conditionally independently given \((H,A)\), from the displayed channels. Put
\[
 g=(1,-2,1)^{\top},
 \qquad
 P(Y=1\mid H=h)=\frac12+\frac18g_h,
\]
using an additional independent uniform exogenous variable. This is a compatible SCM: take \(H_a=H\), \(W_h=W\), \(M_{a,h}=M\), and
\[
 Y_{a,h,m,w}=\mathbf 1\!\left\{U_Y\le\frac12+\frac18g_h\right\}.
\]
These assignments directly verify \(\mathrm{ass{:}treatment\text{-}consistency}\), \(\mathrm{ass{:}z\text{-}consistency}\), \(\mathrm{ass{:}w\text{-}consistency}\), \(\mathrm{ass{:}m\text{-}consistency}\), and \(\mathrm{ass{:}y\text{-}consistency}\). Independence of the exogenous variables verifies \(\mathrm{ass{:}sequential\text{-}treatment}\), \(\mathrm{ass{:}sequential\text{-}outcome}\), \(\mathrm{ass{:}sequential\text{-}mediator}\), and \(\mathrm{ass{:}cross\text{-}world}\), while the construction verifies both proxy exclusions. Every treatment and bridge-denominator probability is positive, and the conditional law of \(H\) given either \(A=1\) or \((M,A=0)\) is uniform and hence strictly positive.

With \(D_\pi=\frac13I_3\), direct multiplication gives
\[
 R_{Z\mid H}D_\pi g=0.
\]
Consequently \(\mathbb E[Y\mid Z,M,A=0]=1/2\). Thus \(h_0(w,m)=1/2\) solves every first-stage bridge equation based on \(Z\), and \(h_1(w)=1/2\) solves every corresponding second-stage bridge equation. In contrast,
\[
 R_{Z^{\star}\mid H}D_\pi g\ne0
\]
because \(R_{Z^{\star}\mid H}\) has full column rank. Moreover, the ratios
\[
 \frac{(R_{Z^{\star}\mid H}D_\pi g)_s}
 {(R_{Z^{\star}\mid H}\pi)_s}
\]
cannot be constant in \(s\). Indeed, if all ratios equalled \(c\), then
\[
 R_{Z^{\star}\mid H}D_\pi(g-c\mathbf 1)=0,
\]
and full column rank would imply \(g=c\mathbf 1\), a contradiction. Bayes' rule therefore shows that \(\mathbb E[Y\mid Z^{\star}=s,M,A=0]\) varies with \(s\). Yet independence of \(W\) gives, for every candidate \(h_0\),
\[
 \sum_{w\in\mathcal X}h_0(w,M)P(W=w\mid Z^{\star}=s,M,A=0)
 =\frac{h_0(0,M)+h_0(1,M)}{2},
\]
which is constant in \(s\) for fixed \(M\). Hence no augmented first-stage bridge exists. This witness preserves the full-rank augmented channel while destroying bridge solvability, proving that the explicit bridge-existence premise is load-bearing and completing the theorem.
\end{proof}
\par\smallskip\noindent \textit{Justification.} The lower bound and constructive attainment certify sharp one-bit restoration of the two completeness clauses; causal identification remains conditional on the other WBC premises and bridge existence. \textit{Gap.} MiaoGengTchetgen2018 gives general finite rank logic but not a positive pair of binary channels restoring both WBC stages simultaneously. \textit{Consumer.} An NLSY97 redesign can add one binary behavior proxy to restore the two completeness clauses, then separately assess the remaining WBC premises and observed bridge existence before invoking identification.

\begin{proposition}[prop:complete-proxy-reduction]\label{prop:complete-proxy-reduction}
For every \(\mathcal S\) satisfying every premise of Wu--Bai--Cui Proposition 3.1 and Theorem 3.1, including their consistency, sequential, cross-world, proxy-exclusion, treatment and mediator positivity, observed-bridge-existence, and completeness conditions, put \(P=P_{\mathcal S}^O\). If \(\mathcal H_{\mathrm{obs}}(P)\neq\varnothing\), Wu--Bai--Cui Proposition 3.1 makes every \(x\in\mathcal H_{\mathrm{obs}}(P)\) latent-valid for \(\mathcal S\), and their Theorem 3.1 gives \(\psi(\mathcal S)=\ell(P)^{\top}x=\mathbb E[h_1(W)\mid A=0]\). Hence \(\ell(P)^{\top}x\) is constant over \(\mathcal H_{\mathrm{obs}}(P)\), \(\mathsf{RI}(P)\) holds, and \(\psi(\mathcal S)=r(P)^{\top}b(P)\), whether or not \(B(P)\) is nonsingular. When \(\mathcal S\in\mathfrak M_{\mathrm{WBC}}^{+}\) also holds, this is exactly the conditional-envelope conclusion.
\end{proposition}
\begin{proof}
Fix \(\mathcal S\) and \(P=P_{\mathcal S}^O\) under the statement's
explicit premises.  The cited
\(\mathrm{lem{:}wbc\text{-}complete\text{-}proxy\text{-}identification}\) says
that the completeness and remaining Wu--Bai--Cui premises lift every solution of
the observed bridge equations to the corresponding latent equations.  Therefore
every \(x\in\mathcal H_{\mathrm{obs}}(P)\) belongs to
\(\mathcal H_{\mathrm{lat}}(\mathcal S)\).  Because the observed fiber is
nonempty, \(b(P)\in\operatorname{col}(B(P))\).  The cited theorem gives, for
every such \(x\),
\[
 \psi(\mathcal S)=\ell(P)^\top x
 =\mathbb E[h_1(W)\mid A=0].
\]
Thus the loading is constant on the whole observed fiber.
\(\mathrm{lem{:}null\text{-}annihilator}\) now implies
\(\ell(P)\in\operatorname{row}(B(P))\), and hence
\(\mathsf{RI}(P)\) holds.  For any \(r(P)\) satisfying
\(B(P)^\top r(P)=\ell(P)\), the same lemma gives
\[
 \psi(\mathcal S)=\ell(P)^\top x=r(P)^\top b(P).
\]
No step uses nonsingularity of \(B(P)\).  If, in addition,
\(\mathcal S\in\mathfrak M_{\mathrm{WBC}}^+\), this is the same equality proved
by \(\mathrm{thm{:}wbc\text{-}target\text{-}span\text{-}identification}\), so
the conditional-envelope result reduces to the cited complete-proxy formula on
the overlap of their assumption sets.
\end{proof}
\par\smallskip\noindent \textit{Justification.} This is the required complete-proxy sanity reduction. \textit{Gap.} WuBaiCui2026 already proves the complete-proxy formula; no novelty is claimed here. \textit{Consumer.} Applied users retain the original WBC formula after successful proxy augmentation.

\begin{lemma}[lem:proximal-id-scope]\label{lem:proximal-id-scope}
\texttt{Shpitser-\allowbreak{}-\allowbreak{}Wood-\allowbreak{}Doughty-\allowbreak{}-\allowbreak{}Tchetgen Tchetgen Assumptions 2.5 and 5.13 impose injectivity at the ordinary and sequential proxy steps,\allowbreak{} and their Theorem 4 identifies the next recursive interventional margin under the stated sequential proxy-\allowbreak{}independence,\allowbreak{} fixing-\allowbreak{}ignorability,\allowbreak{} bridge-\allowbreak{}existence,\allowbreak{} completeness,\allowbreak{} and positivity premises.}
\end{lemma}

\begin{lemma}[lem:proximal-path-specific-scope]\label{lem:proximal-path-specific-scope}
\texttt{Bai-\allowbreak{}-\allowbreak{}Wu-\allowbreak{}-\allowbreak{}Sun-\allowbreak{}-\allowbreak{}Cui study an observed treatment-\allowbreak{}induced recanting witness distinct from a latent baseline confounder;\allowbreak{} their Assumption 6(i)-\allowbreak{}-\allowbreak{}(iii) imposes three stage-\allowbreak{}specific completeness conditions,\allowbreak{} and their Theorem 2.1 uses three bridges to identify a path-\allowbreak{}specific mean on that different graph.}
\end{lemma}

\begin{lemma}[lem:wbc-complete-proxy-identification]\label{lem:wbc-complete-proxy-identification}
Under Wu--Bai--Cui Assumptions 2.1--2.3 and 3.1--3.3, Proposition 3.1 lifts every solution of the two observed outcome-bridge equations to the corresponding latent bridge equations, and Theorem 3.1 identifies the hidden-recanting mean as \(\mathbb E[h_1(W)\mid A=0]\).
\end{lemma}

\begin{lemma}[lem:wbc-por-identification]\label{lem:wbc-por-identification}
In the finite no-baseline-covariate specialization of Wu--Bai--Cui Theorem 3.1, their Assumptions 2.1--2.3 and 3.1--3.3 imply the Proximal Outcome Regression identification formula
\[
 \psi(\mathcal S)=\sum_{w\in\mathcal X}h_1(w)P_{\mathcal S}^O(W=w\mid A=0).
\]
\end{lemma}

\begin{lemma}[lem:duarte-response-program-scope]\label{lem:duarte-response-program-scope}
Duarte--Finkelstein--Knox--Mummolo--Shpitser Algorithm 2 constructs a response-function polynomial program equivalent to the discrete causal bounding problem; their Theorem 1 identifies its minimum and maximum as the sharp lower and upper bounds; and their Algorithm 3 supplies anytime \(\varepsilon\)-sharp optimization.
\end{lemma}
\par\smallskip\noindent \textit{Justification.} This leaf credits the published generic response-function reduction, sharp-bound theorem, and anytime optimization machinery used only for comparison. \textit{Gap.} The cited result is generic; it does not provide the WBC-specific bridge-supported causal-fiber equality, calibrated causal lift, rescue branches, or exact WBC artifacts established here. \textit{Consumer.} It prevents the canonical finite program from being presented as a new generic response-program or optimization result.

\begin{theorem}[thm:proxy-measurable-arm-zero-rescue]\label{thm:proxy-measurable-arm-zero-rescue}
\textbf{Proxy-measurable treatment-zero rescue and failure of universal row-span necessity.} Let \(P\) be an observed law induced by at least one member of \(\mathfrak M_{\mathrm{WBC}}^{+}\). If there is a map \(f:\mathcal X\to\mathcal X\) such that \(P(Y=f(W)\mid A=0)=1\), then every \(\mathcal S\in\mathfrak M_{\mathrm{WBC}}^{+}\) satisfying \(P_{\mathcal S}^O=P\) obeys
\[
\psi(\mathcal S)=\mathbb E_P[f(W)\mid A=0],
\]
whether or not \(\ell(P)\in\operatorname{row}(B(P))\). For a nondegenerate exact witness, replace only the treatment-zero outcome kernel in either SCM of \(\mathrm{def{:}failure\text{-}artifact}\) by \(Y_{0,h,m,w}=w\), retaining its treatment-one outcome kernel. The two modified SCMs remain observationally equivalent members of \(\mathfrak M_{\mathrm{WBC}}^{+}\); their common observed law \(P\) satisfies \(P(Y=1\mid A=0)=5/12\) and \(\ell(P)\notin\operatorname{row}(B(P))\), yet every member of \(\mathfrak M_{\mathrm{WBC}}^{+}\) inducing \(P\) has target \(5/12\). Taking \(f\equiv0\) recovers the deterministic-zero counterexample as a corollary. Thus the universal full-law row-span converse is false.
\end{theorem}
\begin{proof}
Fix an observed law \(P\) satisfying the premise and a map
\(f:\mathcal X\to\mathcal X\).  Let \(\mathcal S\) be any member of
\(\mathfrak M_{\mathrm{WBC}}^{+}\) inducing \(P\).  Define a bridge pair by
\[
 h_0^f(w,m)=f(w)\quad(w,m\in\mathcal X),
 \qquad h_1^f(w)=f(w)\quad(w\in\mathcal X),
\]
and let \(x_f\) be its six-coordinate vector.

The two observable denominator assumptions make every conditional moment in
the coupled bridge system well defined.  Since
\(P(Y=f(W)\mid A=0)=1\), for each \((z,m)\),
\[
\begin{aligned}
 E_P[Y\mid Z=z,M=m,A=0]
 &=E_P[f(W)\mid Z=z,M=m,A=0]\\
 &=\sum_w h_0^f(w,m)P(W=w\mid Z=z,M=m,A=0).
\end{aligned}
\]
For each \(z\),
\[
\begin{aligned}
 &\sum_{w,m}h_0^f(w,m)P(W=w,M=m\mid Z=z,A=1)\\
 &\qquad=\sum_w f(w)P(W=w\mid Z=z,A=1)\\
 &\qquad=\sum_w h_1^f(w)P(W=w\mid Z=z,A=1).
\end{aligned}
\]
By the definition of \(B(P)\) and \(b(P)\), these six equations say
\(B(P)x_f=b(P)\), so \(x_f\in\mathcal H_{\mathrm{obs}}(P)\).

We also verify latent validity rather than infer it from the observed
equations.  If \((h,m)\) is target relevant in the first latent bridge
equation, the latent treatment-zero overlap assumption gives
\(P_{\mathcal S}(H=h,M=m,A=0)>0\).  The observed almost-sure equality
\(Y=f(W)\) on \(\{A=0\}\) therefore remains true after conditioning on
this positive-probability event, and
\[
\begin{aligned}
 E_{P_{\mathcal S}}[Y\mid H=h,M=m,A=0]
 &=E_{P_{\mathcal S}}[f(W)\mid H=h,M=m,A=0]\\
 &=\sum_w h_0^f(w,m)
 P_{\mathcal S}(W=w\mid H=h,M=m,A=0).
\end{aligned}
\]
For every target-relevant \(h\), the second latent bridge equation is the
identity
\[
\begin{aligned}
 &\sum_{w,m}h_0^f(w,m)
 P_{\mathcal S}(W=w,M=m\mid H=h,A=1)\\
 &\qquad=\sum_w f(w)P_{\mathcal S}(W=w\mid H=h,A=1)\\
 &\qquad=\sum_w h_1^f(w)P_{\mathcal S}(W=w\mid H=h,A=1),
\end{aligned}
\]
whose conditionals are defined by the latent treatment-one overlap
assumption.  Hence
\(x_f\in\mathcal H_{\mathrm{lat}}(\mathcal S)\).  Applying
\(\mathrm{lem{:}latent\text{-}bridge\text{-}identifies}\) to this explicit
intersection point gives
\[
 \psi(\mathcal S)=\ell(P)^{\top}x_f
 =\sum_w P(W=w\mid A=0)f(w)
 =E_P[f(W)\mid A=0].
\]
The right side depends only on \(P\), and no row-space premise was used, so
the equality holds for every class member in the observed-law fiber.

For the exact nonconstant witness, modify either member of
\(\mathrm{def{:}failure\text{-}artifact}\) by setting
\(Y_{0,h,m,w}=w\) and leave every other response unchanged.  The consistency,
sequential, cross-world, exclusion, denominator, and overlap arguments in the
counterpair proof are unchanged: in particular, the new treatment-zero
potential outcome uses only \(U_W\) and is still independent of the mediator
shock \(U_M\).  The bridge pair
\[
 h_0(w,m)=w,\qquad h_1(w)=w
\]
is both observed-valid and latent-valid by the preceding calculation, so the
modified models remain in \(\mathfrak M_{\mathrm{WBC}}^{+}\).  The models
remain observationally equivalent because their treatment-zero mechanisms
are modified identically, while their only difference is still the
treatment-one mediator kernel and the treatment-one outcome remains
Bernoulli with mean \(1/2\).  The law of \((A,Z,W,M)\), and hence both
\(B(P)\) and \(\ell(P)\), is unchanged.  Thus the same null vector
\((0,0,0,0,1,-1)^{\top}\) proves
\(\ell(P)\notin\operatorname{row}(B(P))\).  On the other hand,
\[
 P(Y=1\mid A=0)=P(W=1\mid A=0)
 =\sum_h\pi _0(h)q_h=\frac5{12}.
\]
The already-proved general part, with \(f(w)=w\), therefore fixes the target
at \(5/12\) for every class member inducing this law.  Taking \(f\equiv0\)
gives the stated deterministic-zero corollary and proves that row span is not
universally necessary.
\end{proof}
\par\smallskip\noindent \textit{Justification.} This is a constructive sufficient identification branch: an observable treatment-zero outcome that is a deterministic function of the outcome-inducing proxy identifies the hidden-recanting target even when the loading fails the bridge row-space test. \textit{Gap.} The failed-row-span witness limits only a universal necessity/converse within this note's WBC envelope. It is not a refutation of a named published estimator or workflow, and it does not claim that proxy measurability exhausts all rescue laws. \textit{Consumer.} For the NLSY97 high-school-track-to-wage interpretation, a general-track wage indicator determined by the measured outcome proxy supplies a constructive identification diagnostic even under incomplete proxy rank.

\begin{theorem}[thm:canonical-full-law-fiber-characterization]\label{thm:canonical-full-law-fiber-characterization}
\textbf{Exact canonical-fiber characterization of the maximal rescue class and of the positive-interior obstruction.}
For a finite set \(E\), write
\[
 \Delta(E)=\left\{q\in[0,1]^E:\sum_{e\in E}q(e)=1\right\}.
\]
Define the finite response-type sets
\[
\begin{aligned}
 \mathcal R_A&=\mathcal A,&
 \mathcal R_H&=\mathcal H^{\mathcal A},&
 \mathcal R_Z&=\mathcal X^{\mathcal A\times\mathcal H},\\
 \mathcal R_W&=\mathcal X^{\mathcal H},&
 \mathcal R_M&=\mathcal X^{\mathcal A\times\mathcal H},&
 \mathcal R_Y&=\mathcal X^{\mathcal A\times\mathcal H\times\mathcal X\times\mathcal X},
\end{aligned}
\]
and \(\mathcal R_{HY}=\mathcal R_H\times\mathcal R_Y\).  A canonical parameter is
\[
 \theta=(\theta_A,\theta_{HY},\theta_Z,\theta_W,\theta_M)
 \in\Delta(\mathcal R_A)\times\Delta(\mathcal R_{HY})
 \times\Delta(\mathcal R_Z)\times\Delta(\mathcal R_W)
 \times\Delta(\mathcal R_M).
\]
For \(r=(r_A,r_H,r_Y,r_Z,r_W,r_M)\), put
\[
 q_\theta(r)=\theta_A(r_A)\theta_{HY}(r_H,r_Y)
 \theta_Z(r_Z)\theta_W(r_W)\theta_M(r_M)
\]
and recursively set
\[
 a_r=r_A,\quad h_r=r_H(a_r),\quad z_r=r_Z(a_r,h_r),\quad
 w_r=r_W(h_r),\quad m_r=r_M(a_r,h_r),\quad
 y_r=r_Y(a_r,h_r,m_r,w_r).
\]
For a predicate \(C\) on the finite type space and a real function \(g\), let
\(\mu_\theta[g\mathbf 1_C]=\sum_rq_\theta(r)g(r)\mathbf 1_C(r)\), and define
\[
 p_\theta(a,z,w,m,y)=
 \mu_\theta[\mathbf 1\{(a_r,z_r,w_r,m_r,y_r)=(a,z,w,m,y)\}].
\]
Writing \(h_r^0=r_H(0)\), define the canonical causal objective
\[
 T(\theta)=\sum_rq_\theta(r)
 r_Y\!\left(0,h_r^0,r_M(1,h_r^0),r_W(h_r^0)\right).
 \tag{12}
\]

For an observed law \(P\), let \(\mathcal F(P)\) be the set of pairs
\((\theta,x)\), where \(\theta\) is as above and
\[
 x=(h_0(0,0),h_0(1,0),h_0(0,1),h_0(1,1),h_1(0),h_1(1))^\top\in\mathbb R^6,
\]
satisfying all of the following explicit constraints.

First, \(p_\theta(o)=P(o)\) for every
\(o\in\mathcal A\times\mathcal X^4\), and
\[
 \sum_{w,y}P(0,z,w,m,y)>0\quad(z,m\in\mathcal X),
 \qquad
 \sum_{w,m,y}P(1,z,w,m,y)>0\quad(z\in\mathcal X).
 \tag{13}
\]
Second, with
\[
\begin{aligned}
 u_h&=\mu_\theta[\mathbf 1\{r_H(0)=h\}],&
 v_{hm}&=\mu_\theta[\mathbf 1\{r_M(1,h)=m\}],\\
 j^0_{hm}&=\mu_\theta[\mathbf 1\{a_r=0,h_r=h,m_r=m\}],&
 j^1_h&=\mu_\theta[\mathbf 1\{a_r=1,h_r=h\}],
\end{aligned}
\]
the target-relevant overlap constraints are
\[
 u_hv_{hm}>0\Longrightarrow j^0_{hm}>0,
 \qquad
 u_h>0\Longrightarrow j^1_h>0.
 \tag{14}
\]
Third, \(B(P)x=b(P)\).  Fourth, on defining
\[
 C^0_{hm}=\{r:a_r=0,h_r=h,m_r=m\},
 \qquad C^1_h=\{r:a_r=1,h_r=h\},
\]
the cross-multiplied latent bridge equations are
\[
 \mu_\theta[y_r\mathbf 1_{C^0_{hm}}]
 =\sum_{w\in\mathcal X}h_0(w,m)
   \mu_\theta[\mathbf 1_{C^0_{hm}\cap\{w_r=w\}}]
 \quad\text{whenever }u_hv_{hm}>0,
 \tag{15}
\]
and
\[
 \sum_{w,m\in\mathcal X}h_0(w,m)
   \mu_\theta[\mathbf 1_{C^1_h\cap\{w_r=w,m_r=m\}}]
 =\sum_{w\in\mathcal X}h_1(w)
   \mu_\theta[\mathbf 1_{C^1_h\cap\{w_r=w\}}]
 \quad\text{whenever }u_h>0.
 \tag{16}
\]

Then, for every \(P\) induced by \(\mathfrak M_{\mathrm{WBC}}^{+}\),
\[
 \{\psi(\mathcal S):\mathcal S\in\mathfrak M_{\mathrm{WBC}}^{+},\ 
 P_{\mathcal S}^O=P\}
 =\Gamma(P):=\{T(\theta):(\theta,x)\in\mathcal F(P)\}.
 \tag{17}
\]
In particular, with
\[
 \underline\psi(P)=\inf\Gamma(P),\qquad
 \overline\psi(P)=\sup\Gamma(P),
 \tag{18}
\]
the unique maximal observable subclass on which the causal target is constant on
every full-law fiber is
\[
 \mathcal I_{\max}=
 \{P:\mathcal F(P)\ne\varnothing,\ 
       \underline\psi(P)=\overline\psi(P)\}.
 \tag{19}
\]
For every induced \(P\), two compatible class members with different targets
exist if and only if
\(\underline\psi(P)<\overline\psi(P)\).  A pair of feasible points with unequal
values of \(T\) is therefore a finite polynomial nonidentification certificate;
after the positive denominators in (13) are cleared, all its equalities and
inequalities are checked directly from the displayed finite sums.

The row-span and proxy-measurable rescue subclasses are contained in
\(\mathcal I_{\max}\).  On the other hand, the strictly positive law
\(P^{\dagger}\) satisfies
\[
 \ell(P^{\dagger})\notin\operatorname{row}(B(P^{\dagger})),
 \qquad
 \overline\psi(P^{\dagger})-\underline\psi(P^{\dagger})\ge\frac1{24}.
 \tag{20}
\]
Finally, define the explicit positive-interior obstruction set
\[
\begin{aligned}
 \mathcal E_{\mathrm{int}}=\{P:\;&\mathcal F(P)\ne\varnothing,
 \ P(o)>0\ \text{for every }o\in\mathcal A\times\mathcal X^4,\\
 &\exists v\in\mathbb R^6\ [B(P)v=0,\ \ell(P)^\top v\ne0],\\
 &T(\theta_1)=T(\theta_2)\ \text{for every }
   (\theta_1,x_1),(\theta_2,x_2)\in\mathcal F(P)\}.
\end{aligned}
 \tag{21}
\]
The proposed universal positive-interior converse is true exactly when
\(\mathcal E_{\mathrm{int}}=\varnothing\).  Thus (21), rather than observed-law
positivity or failed row span alone, is the exact remaining obstruction.  This
theorem gives the maximal-class characterization and an exactly checkable
two-law certificate; it does not assert, without a further latent-tilt argument,
that \(\mathcal E_{\mathrm{int}}\) is empty.
\end{theorem}
\begin{proof}
For an observed law \(P\), denote its full causal fiber by
\[
 \mathcal C(P)=\{\mathcal S\in\mathfrak M_{\mathrm{WBC}}^{+}:
 P_{\mathcal S}^O=P\}.
\]
We prove (17) by constructing mutually inverse representations at the level
of induced full-law target values.

\emph{From a class member to a feasible canonical point.}
Fix \(\mathcal S\in\mathcal C(P)\).  Because all endogenous alphabets are
finite, every realization of an exogenous shock determines a deterministic
response table.  Let \(R_A\) be the resulting treatment response, let
\((R_H,R_Y)\) be the joint response table for the district \(\{H,Y\}\),
and let \(R_Z,R_W,R_M\) be the other three response tables.  Compatibility
with \(\mathcal G_{\mathrm{WBC}}\) means that exogenous blocks belonging to
distinct districts are independent, while arbitrary dependence is allowed
inside the \(\{H,Y\}\) block.  Therefore their pushforward laws have the
form
\[
 \theta_A\otimes\theta_{HY}\otimes\theta_Z\otimes\theta_W
 \otimes\theta_M,
\]
on exactly the finite response-type spaces displayed in the theorem.  In
particular, the joint mass of a response tuple \(r\) is \(q_\theta(r)\).

Consistency and recursive structural substitution give
\[
 A=r_A,\quad H=r_H(r_A),\quad Z=r_Z(A,H),\quad W=r_W(H),\quad
 M=r_M(A,H),\quad Y=r_Y(A,H,M,W).
\]
These are precisely the recursions defining
\((a_r,h_r,z_r,w_r,m_r,y_r)\).  It follows cell by cell that
\[
 p_\theta(a,z,w,m,y)=P_{\mathcal S}^O(a,z,w,m,y)=P(a,z,w,m,y).
 \tag{22}
\]
The two denominator inequalities in (13) follow from the corresponding
member assumptions.

Independence of \(R_M\) from the other response blocks and consistency give
\[
\begin{aligned}
 u_h&=P_{\mathcal S}(H_0=h),&
 v_{hm}&=P_{\mathcal S}(M_{1,h}=m),\\
 j^0_{hm}&=P_{\mathcal S}(A=0,H=h,M=m),&
 j^1_h&=P_{\mathcal S}(A=1,H=h).
\end{aligned}
 \tag{23}
\]
Thus the two latent transport-overlap member assumptions imply (14).
Latent-valid observed bridge existence supplies an
\[
 x\in\mathcal H_{\mathrm{lat}}(\mathcal S)
      \cap\mathcal H_{\mathrm{obs}}(P).
\]
Observed validity is exactly \(B(P)x=b(P)\).  For every target-relevant
\((h,m)\), multiply the first conditional latent bridge equation by the
positive number \(j^0_{hm}\) in (23).  The result is (15), because
\(C^0_{hm}=\{A=0,H=h,M=m\}\) under the response recursion.  Similarly,
multiplying the second conditional latent bridge equation by the positive
number \(j^1_h\) gives (16).  Hence \((\theta,x)\in\mathcal F(P)\).

Finally, the same response table \(r_H(0)\) is used in both occurrences of
the hidden witness in
\[
 r_Y\!\left(0,r_H(0),r_M(1,r_H(0)),r_W(r_H(0))\right).
\]
Taking its expectation under \(q_\theta\) therefore gives, by the definition
of the nested potential outcome,
\[
 T(\theta)=E_{\mathcal S}
 [Y_{0,H_0,M_{1,H_0},W_{H_0}}]=\psi(\mathcal S).
 \tag{24}
\]
Equations (22)--(24) prove the inclusion of the left side of (17) in
\(\Gamma(P)\).

\emph{From a feasible canonical point to a class member.}
Conversely, fix \((\theta,x)\in\mathcal F(P)\).  On the finite product
response space, draw
\[
 R_A,\ (R_H,R_Y),\ R_Z,\ R_W,\ R_M
\]
independently across the five displayed blocks with marginal laws
\(\theta_A,\theta_{HY},\theta_Z,\theta_W,\theta_M\), respectively.  Define
the recursive SCM
\[
 A=R_A,\quad H=R_H(A),\quad Z=R_Z(A,H),\quad W=R_W(H),\quad
 M=R_M(A,H),\quad Y=R_Y(A,H,M,W).
 \tag{25}
\]
Only \(R_H\) and \(R_Y\) may be dependent, exactly realizing the district
\(\{H,Y\}\); all directed arguments in (25) are allowed parents in
\(\mathcal G_{\mathrm{WBC}}\).  Thus the constructed SCM, denoted
\(\mathcal S_\theta\), is compatible with that graph, and all five
consistency assumptions hold identically.

For completeness, we verify rather than assume the remaining structural
member properties.  The random variable \(A\) is a function of \(R_A\),
whereas \(\{Y_{a,H_a,m,W_{H_a}},H_a\}\) is a function of
\((R_H,R_Y,R_W)\), proving sequential treatment independence.  Each
\(M_{a,h}=R_M(a,h)\) is independent of \((H,A)\), and it is also independent
of \(\{Y_{a,H_a,m,W_{H_a}},H_a\}\), proving mediator independence and the
cross-world condition.  Conditional on \((H,A)\), the event conditioned on
depends only on \((R_H,R_A)\); the realized \(M\) then uses the independent
block \(R_M\), while \(Y_{A,H,m,W_H}\) uses \((R_Y,R_W)\).  This proves
the sequential-outcome condition.  In the same way, conditional on
\((H,A)\), \(Z\) uses the independent block \(R_Z\), so it is independent
of \((Y,M)\).  Conditional on \(H\), \(W\) uses the independent block
\(R_W\), whereas \((A,Z,M)\) uses \((R_A,R_H,R_Z,R_M)\); hence the two
proxy-exclusion conditions hold.  These conditional factorizations remain
valid on every conditioning event of positive probability, which is the
finite-law meaning of the asserted conditional independences.

By construction, the observed cell probabilities of (25) are
\(p_\theta(o)\), so the first feasibility constraint yields
\(P_{\mathcal S_\theta}^O=P\).  The strict inequalities in (13) give the two
observable denominator assumptions.  The equalities in (23) also hold for
the constructed SCM.  Therefore (14) gives both target-relevant latent
transport-overlap assumptions.

It remains only to prove bridge-supported membership.  Feasibility already
gives \(B(P)x=b(P)\), hence
\(x\in\mathcal H_{\mathrm{obs}}(P)\).  If \(u_hv_{hm}>0\), (14) gives
\(j^0_{hm}>0\).  Dividing (15) by this positive number converts its left side
to
\(E_{P_{\mathcal S_\theta}}[Y\mid H=h,M=m,A=0]\) and its right side to
\[
 \sum_w h_0(w,m)
 P_{\mathcal S_\theta}(W=w\mid H=h,M=m,A=0).
\]
This is exactly the first latent bridge equation.  If \(u_h>0\), (14) gives
\(j^1_h>0\); dividing (16) by \(j^1_h\) gives exactly the second latent
bridge equation.  Hence
\[
 x\in\mathcal H_{\mathrm{lat}}(\mathcal S_\theta)
      \cap\mathcal H_{\mathrm{obs}}(P),
\]
so the latent-valid bridge-existence assumption holds.  We have proved
\(\mathcal S_\theta\in\mathcal C(P)\).  Direct substitution in (25), using
one and the same coordinate \(R_H(0)\) throughout the nested counterfactual,
again gives
\[
 \psi(\mathcal S_\theta)=T(\theta).
 \tag{26}
\]
This proves the reverse inclusion in (17), and hence equality.

We now derive each characterization.  If \(P\) is induced by the class, its
fiber and therefore \(\Gamma(P)\) are nonempty.  Also
\(T(\theta)\in[0,1]\), since it is the expectation of a binary response.
For any nonempty subset \(G\subseteq\mathbb R\), every element of \(G\) is
equal if and only if \(\inf G=\sup G\).  Applying this elementary fact to
the exact equality (17) proves that the target is constant on the full-law
fiber precisely when
\(\underline\psi(P)=\overline\psi(P)\).  Conversely, if
\(\mathcal F(P)\ne\varnothing\), the construction (25) shows that \(P\) is
induced by a member of the class.  Thus (19) contains every and only observed
law whose full-law fiber identifies the target.  Any observable subclass with
that property is a subset of (19), which proves uniqueness and maximality.

The same elementary order argument shows
\[
 \underline\psi(P)<\overline\psi(P)
 \quad\Longleftrightarrow\quad
 \exists (\theta_1,x_1),(\theta_2,x_2)\in\mathcal F(P)
 \text{ with }T(\theta_1)\ne T(\theta_2).
\]
By (25)--(26), the two feasible points construct two compatible class members
with the same observed law and unequal targets.  Conversely, any such two
members give two such feasible points by the forward construction.  This is
the asserted exact two-SCM certificate.  It is finite: simplex membership,
the response sums, (14)--(16), and the two unequal objective values comprise
finitely many polynomial equalities, weak inequalities, strict inequalities,
and finite disjunctions.  The entries of \(B(P)\) and \(b(P)\) are rational
functions of observed cells whose denominators are positive by (13);
multiplication by those denominators turns their bridge equations into
polynomial equalities.  Thus an exhibited pair is checked directly without
any endpoint-attainment premise.

If \(\ell(P)\in\operatorname{row}(B(P))\),
\(\mathrm{thm{:}wbc\text{-}target\text{-}span\text{-}identification}\)
shows that every member of \(\mathcal C(P)\) has the same target.  If instead
the proxy-measurable arm-zero premise holds,
\(\mathrm{thm{:}proxy\text{-}measurable\text{-}arm\text{-}zero\text{-}rescue}\)
gives the same conclusion.  The endpoint characterization just proved puts
both sufficient subclasses inside \(\mathcal I_{\max}\).

By \(\mathrm{prop{:}scm\text{-}counterpair}\), \(P^{\dagger}\) is strictly
positive, its loading is not in the row space, and its full-law fiber contains
target values \(19/48\) and \(7/16\).  Exactness of (17) therefore gives
\[
 \overline\psi(P^{\dagger})-\underline\psi(P^{\dagger})
 \ge \frac7{16}-\frac{19}{48}=\frac1{24},
\]
which proves (20).

Finally, consider the universal positive-interior converse: at every induced
law with all observed cells positive, failure of row span should imply two
compatible class members with unequal targets.  Whenever
\(\mathcal F(P)\ne\varnothing\), feasibility supplies an \(x\) with
\(B(P)x=b(P)\), so \(b(P)\in\operatorname{col}(B(P))\).  By
\(\mathrm{lem{:}null\text{-}annihilator}\), row-span failure is then
equivalent to the existence of \(v\) with
\(B(P)v=0\) and \(\ell(P)^{\top}v\ne0\).  By (17), absence of two unequal
compatible targets is equivalent to constancy of \(T\) on
\(\mathcal F(P)\).  Consequently a counterexample to the universal converse
is exactly an element of \(\mathcal E_{\mathrm{int}}\).  The converse is
therefore true if and only if \(\mathcal E_{\mathrm{int}}=\varnothing\).
Nothing in this equivalence proves that emptiness; it isolates it as the
remaining latent-tilt question, exactly as claimed.
\end{proof}
\par\smallskip\noindent \textit{Justification.} The theorem proves the WBC-specific equality between the compatible-SCM causal fiber and the finite canonical feasible-image \(\Gamma(P)\); endpoint equality then characterizes the maximal point-identification subclass, and unequal feasible objectives give exact finite nonidentification certificates. \textit{Gap.} No accepted-bank in-repository comparator is asserted for this theorem. Its scope is the exact bridge-supported hidden-witness WBC causal-fiber equality, the induced maximal-class characterization, and the stated finite artifacts; it does not claim endpoint attainment, a terminating solver, a branch-and-bound or sum-of-squares certificate theorem, certified approximation, or inference. The sole residual is emptiness of \(\mathcal E_{\mathrm{int}}\). \textit{Consumer.} An analyst can form the exact population feasible set, test point identification through equality of the two infimum/supremum definitions, or provide two exactly checkable feasible points with unequal objectives as a nonidentification certificate; no terminating numerical endpoint computation is promised.

\begin{openendedquestion}[oeq:positive-interior-converse]\label{oeq:positive-interior-converse}
Is the positive-interior obstruction set displayed in \(\mathrm{thm{:}canonical\text{-}full\text{-}law\text{-}fiber\text{-}characterization}\) empty?  Equivalently, for every observed law \(P\) induced by \(\mathfrak M_{\mathrm{WBC}}^{+}\) that assigns positive probability to every element of \(\mathcal A\times\mathcal X^4\) and satisfies \(\ell(P)\notin\operatorname{row}(B(P))\), do there exist \(\mathcal S_1,\mathcal S_2\in\mathfrak M_{\mathrm{WBC}}^{+}\) such that \(P_{\mathcal S_1}^O=P_{\mathcal S_2}^O=P\) and \(\psi(\mathcal S_1)\ne\psi(\mathcal S_2)\)?
\end{openendedquestion}
Fix \(P\), a feasible point \((\theta,x)\in\mathcal F(P)\), a vector \(v=(v_0,v_1)\in\mathbb R^6\), and \(t\ne0\). Suppose, for contradiction, that \((\theta,x+t v)\in\mathcal F(P)\) as well. Write
\[
 \pi_0(h)=\Pr_\theta\{R_H(0)=h\},\qquad
 q_h(w)=\Pr_\theta\{R_W(h)=w\},\qquad
 s_{1h}(m)=\Pr_\theta\{R_M(1,h)=m\}.
\]
For every \(h\) with \(\pi_0(h)>0\) and every \(m\) with \(s_{1h}(m)>0\), the overlap implication (14) gives \(j^0_{hm}>0\). Subtracting equation (15) for \(x\) from the same equation for \(x+t v\), and using \(t\ne0\), yields
\[
 0=\sum_w v_0(w,m)\,
 \mu_\theta[\mathbf 1_{C^0_{hm}\cap\{w_r=w\}}].
\]
The product factorization of \(q_\theta\) makes \(R_W\) independent of \((R_A,R_H,R_M)\), so the last display equals
\[
 j^0_{hm}\sum_w q_h(w)v_0(w,m).
\]
Division by \(j^0_{hm}>0\) therefore gives
\[
 \sum_wq_h(w)v_0(w,m)=0.
 \tag{A}
\]
For every \(h\) with \(\pi_0(h)>0\), (14) also gives \(j^1_h>0\). Subtracting equation (16) at the two bridge vectors and again using the product factorization gives
\[
 \sum_m s_{1h}(m)\sum_wq_h(w)v_0(w,m)
 =\sum_wq_h(w)v_1(w).
\]
Every positive-weight term on the left vanishes by (A), and every remaining term has \(s_{1h}(m)=0\). Hence
\[
 \sum_wq_h(w)v_1(w)=0
 \quad\text{for every }h\text{ with }\pi_0(h)>0.
 \tag{B}
\]
Finally, the canonical product factorization and recursive evaluation imply
\[
 P(W=w\mid A=0)=\sum_h\pi_0(h)q_h(w).
\]
Because \(\ell(P)\) has zero \(h_0\)-coordinates and its two \(h_1\)-coordinates are \(P(W=w\mid A=0)\), equation (B) gives
\[
 \ell(P)^{\top}v
 =\sum_h\pi_0(h)\sum_wq_h(w)v_1(w)=0.
\]
Thus a nonzero-loading observed null direction cannot be lifted by changing the bridge alone at a fixed canonical response law. Any affirmative proof must construct a simultaneous, factor-preserving change in \(\theta\); the frozen assumptions do not provide the latent slack or rank condition needed to do so. The already established canonical-fiber theorem still reduces either resolution to a finite quantified polynomial certificate, but it does not decide which certificate exists.
\par\smallskip\noindent \textit{Justification.} The exact maximal-class theorem reduces the remaining universal positive-interior converse to one explicit emptiness question. \textit{Gap.} Observed positivity and failed bridge row span alone do not construct the required latent null tilt; a compatible two-law perturbation is still needed. \textit{Consumer.} Resolving this question determines whether every strictly positive failed-row-span law is causally nonidentified within the fixed WBC class.

\section{Interpretation}

For the NLSY97 high-school-track-to-wage application, \(\Gamma(P)\) is the entire range of the same-unit wage-path mean compatible with the fixed proxy-augmented SCM and observed multinomial law. Equality of its endpoints certifies point identification; unequal feasible objectives exhibit two concrete compatible SCMs and should be reported as partial rather than point identification. The calibrated rank-two benchmark shows that binary coarsening need not destroy scalar identification, while the positive counterpair provides a concrete failed-row-span law with genuine causal ambiguity. The proxy-measurable theorem is a separate constructive rescue branch, not a critique of a specified estimation workflow.

\section*{Technical internal limitation (diagnostic only)}

The calibrated-subclass theorem and the broad conditional envelope are sufficient identification results. The exact canonical response-type program characterizes the full causal fiber, the sharp infimum and supremum, and the maximal point-identification subclass, while the explicit failure pair proves nonidentification at one positive law. What remains open is the positive-interior converse: whether every strictly positive induced law with failed loading row span has unequal causal-fiber endpoints, equivalently whether \(\mathcal E_{\mathrm{int}}\) is empty. The population characterization does not prove endpoint attainment, a terminating endpoint solver, a branch-and-bound or sum-of-squares certificate theorem, certified approximation, or root-\(n\)/efficiency theory.

\section{Honest scope}

Null-annihilator algebra, generalized inverses, finite rank counting, ordinary completeness, and the generic response-function polynomial reduction and optimization machinery of Duarte et al. are not contributions. The broad theorem is only a transparent conditional envelope: membership assumes \(\mathcal H_{\mathrm{lat}}(\mathcal S)\cap\mathcal H_{\mathrm{obs}}(P_{\mathcal S}^O)\neq\varnothing\), which automatically implies observed bridge solvability, together with two novel target-relevant latent transport-overlap restrictions that are not testable from \(O\). The calibrated subclass is the strongest honest proved sufficient subclass currently available. It derives latent-valid observed bridge realizability directly from calibrated binary measurement, affine treatment-one mediator risk, and additive treatment-zero outcome means, but it conditions on rather than derives the two latent transport restrictions. The exact success SCM lies in this primitive subclass and identifies a positive law where both three-state completeness conditions fail; no claim is made that the full calibrated assumption package is logically weaker than the complete Wu--Bai--Cui model or the Bai--Wu--Sun--Cui package. Bai--Wu--Sun--Cui permit nonunique bridge solutions, so the comparison concerns their observed-recanting-witness graph, distinct latent confounder, target, and completeness-based causal lift rather than bridge uniqueness. Neither theorem subsumes the other: this proposal does not cover their graph, three-bridge target, four identification strategies, or inference, and their result does not cover this hidden-witness same-unit target or the exact rank-deficient success member. The exact canonical theorem identifies the full WBC causal-fiber range \(\Gamma(P)\) and the unique maximal observable point-identification subclass by equality of its infimum and supremum; it gives an exact population characterization and checkable two-feasible-point certificates, not a new generic reduction, endpoint-attainment theorem, terminating solver, branch-and-bound or sum-of-squares guarantee, certified approximation, or inference procedure. The completed failure artifact fixes the same-unit joint response of \(H_0\) and \(H_1\) without changing any kernel or target value. The remaining universal question is only whether the explicit positive-interior obstruction set \(\mathcal E_{\mathrm{int}}\) is empty. The complete-proxy case reduces to Wu--Bai--Cui Theorem 3.1. The augmentation result is input-specific to the displayed rank-two \(Z\)-channel and proves only sharp one-bit restoration of the two completeness clauses. Full augmented-channel rank establishes neither the remaining causal, consistency, exclusion, support or positivity premises nor observed bridge existence; Wu--Bai--Cui identification follows only when those additional premises hold. No positivity of individual \(W\) or \(Y\) values and no root-\(n\) or efficiency result are imposed or claimed.

\begin{thebibliography}{99}

  \bibitem{WuBaiCui2026} Sihan Wu, Yang Bai, and Yifan Cui (2026). Proximal Mediation Analysis with Hidden Recanting Witnesses. arXiv:2606.17600v1, especially Assumptions 2.1--2.3 and 3.1--3.3, Proposition 3.1, and Theorems 3.1 and 4.1.

  \bibitem{BaiWuSunCui2026} Yang Bai, Sihan Wu, Baoluo Sun, and Yifan Cui (2026). Proximal Path-Specific Inference. arXiv:2605.09462v1, 10 May 2026, especially Section 2 and Figure 1(c), Assumption 5, Assumption 6(i)--(iii), and Theorem 2.1, equations (4)--(7).

  \bibitem{ZhangLiMiaoTchetgen2023} Jeffrey Zhang, Wei Li, Wang Miao, and Eric Tchetgen Tchetgen (2023). Proximal Causal Inference without Uniqueness Assumptions. Statistics and Probability Letters 198, 109836, especially Propositions 2.3 and 2.5 and Lemmas 2.4 and 2.6.

  \bibitem{BennettEtAl2023} Andrew Bennett, Nathan Kallus, Xiaojie Mao, Whitney Newey, Vasilis Syrgkanis, and Masatoshi Uehara (2023). Inference on Strongly Identified Functionals of Weakly Identified Functions. Proceedings of Machine Learning Research 195.

  \bibitem{KallusMaoUehara2021} Nathan Kallus, Xiaojie Mao, and Masatoshi Uehara (2021 manuscript; 2024 journal version). Causal Inference Under Unmeasured Confounding with Negative Controls: A Minimax Learning Approach. Journal of Machine Learning Research 25.

  \bibitem{Santos2011} Andres Santos (2011). Instrumental Variable Methods for Recovering Continuous Linear Functionals. Journal of Econometrics 161(2), 129--146.

  \bibitem{SeveriniTripathi2012} Thomas A. Severini and Gautam Tripathi (2012). Efficiency Bounds for Estimating Linear Functionals of Nonparametric Regression Models with Endogenous Regressors. Journal of Econometrics 170(2), 491--498.

  \bibitem{MiaoGengTchetgen2018} Wang Miao, Zhi Geng, and Eric Tchetgen Tchetgen (2018). Identifying Causal Effects with Proxy Variables of an Unmeasured Confounder. Biometrika 105(4), 987--993.

  \bibitem{GhassamiZhangShpitserTchetgen2026} AmirEmad Ghassami, Yunshu Zhang, Ilya Shpitser, and Eric Tchetgen Tchetgen (2026). Partial Identification of Causal Effects Using Proxy Variables. arXiv:2304.04374v4, 5 August 2026.

  \bibitem{DukesShpitserTchetgen2023} Oliver Dukes, Ilya Shpitser, and Eric Tchetgen Tchetgen (2023). Proximal Mediation Analysis. Biometrika 110(4), 973--987.

  \bibitem{ShpitserWoodDoughtyTchetgen2023} Ilya Shpitser, Zach Wood-Doughty, and Eric J. Tchetgen Tchetgen (2023). The Proximal ID Algorithm. Journal of Machine Learning Research 24(188), 1--46, especially Assumptions 2.5 and 5.13, Theorem 4, and Section 5.3.

  \bibitem{RepoProximalPhase} CausalSmith record \(\mathrm{eid\_proximal\_phase\_v1}\): finite null-annihilator geometry and bridge-selector diagnostics.

  \bibitem{RepoProximalCompleteness} CausalSmith record \(\mathrm{eid\_proximal\_completeness\_v1}\): an unobservable operator threshold with an incomplete causal link.

  \bibitem{RepoProximalDID} CausalSmith record \(\mathrm{eid\_proximal\_did\_bridge\_frontier\_v1}\): known null-annihilator algebra without a full-law converse.

  \bibitem{RepoProxyTargetSpan} CausalSmith proposal \(\mathrm{scm\_proxy\_targetspan\_transport}\): target-span identification and a full-law converse for a transport SCM.

  \bibitem{DuarteFinkelsteinKnoxMummoloShpitser2024} Guilherme Duarte, Noam Finkelstein, Dean Knox, Jonathan Mummolo, and Ilya Shpitser (2024). An Automated Approach to Causal Inference in Discrete Settings. Journal of the American Statistical Association 119, 1778--1793.

\end{thebibliography}

\end{document}